\documentclass[11pt,a4paper]{report}
\usepackage[margin=1in,top=1.1in,bottom=1.1in]{geometry}
\usepackage[utf8]{inputenc}
\usepackage[T1]{fontenc}
\usepackage{lmodern}
\usepackage{listings}
\usepackage{xcolor}
\usepackage{amsmath,amssymb}
\usepackage{booktabs}
\usepackage{longtable}
\usepackage{array}
\usepackage[colorlinks=true,linkcolor=blue!60!black,urlcolor=blue!70!black]{hyperref}
\usepackage[most,breakable]{tcolorbox}
\usepackage{fancyhdr}
\usepackage{titlesec}
\usepackage{enumitem}
\usepackage{microtype}
\usepackage{multicol}
\usepackage{mdframed}

% ── Colours ──────────────────────────────────────────────────
\definecolor{codebg}{RGB}{250,250,252}
\definecolor{pyblue}{RGB}{0,0,188}
\definecolor{pygreen}{RGB}{0,128,0}
\definecolor{pyorange}{RGB}{186,74,0}
\definecolor{pygray}{RGB}{120,120,120}
\definecolor{accentblue}{RGB}{0,82,165}
\definecolor{noteblue}{RGB}{232,244,255}
\definecolor{noteborder}{RGB}{0,100,200}
\definecolor{compyellow}{RGB}{255,251,222}
\definecolor{compborder}{RGB}{190,140,0}
\definecolor{tipgreen}{RGB}{232,255,232}
\definecolor{tipborder}{RGB}{0,140,0}
\definecolor{warnred}{RGB}{255,235,235}
\definecolor{warnborder}{RGB}{190,0,0}

% ── Listing Styles ────────────────────────────────────────────
\lstdefinestyle{py}{
  language=Python,
  backgroundcolor=\color{codebg},
  basicstyle=\ttfamily\small,
  keywordstyle=\color{pyblue}\bfseries,
  stringstyle=\color{pyorange},
  commentstyle=\color{pygreen}\itshape,
  numberstyle=\tiny\color{pygray},
  numbers=left, stepnumber=1, numbersep=8pt,
  frame=single, framerule=0.4pt, rulecolor=\color{gray!40},
  breaklines=true, tabsize=4, showstringspaces=false, keepspaces=true,
  xleftmargin=18pt, xrightmargin=4pt, aboveskip=8pt, belowskip=8pt,
  morekeywords={None,True,False,self,cls,print,range,len,int,str,list,
    dict,set,tuple,float,bool,type,super,enumerate,zip,map,filter,
    sorted,reversed,sum,min,max,abs,isinstance,hasattr,getattr,
    setattr,append,extend,pop,insert,remove,index,count,values,keys,
    items,get,update,clear,copy,defaultdict,Counter,deque,heapq,
    bisect,lru_cache,cache,dataclass,property,staticmethod,
    classmethod,yield,lambda,assert,pass,raise,del,with,as,
    from,import,global,nonlocal},
}
\lstdefinestyle{cpp}{
  language=C++,
  backgroundcolor=\color{codebg},
  basicstyle=\ttfamily\small,
  keywordstyle=\color{pyblue}\bfseries,
  stringstyle=\color{pyorange},
  commentstyle=\color{pygreen}\itshape,
  numberstyle=\tiny\color{pygray},
  numbers=left, frame=single, breaklines=true, tabsize=4,
  showstringspaces=false, xleftmargin=18pt, aboveskip=8pt, belowskip=8pt,
}
\lstset{style=py}

% ── tcolorbox environments ────────────────────────────────────
\newtcolorbox{notebox}[1][Note]{
  colback=noteblue, colframe=noteborder,
  title=\textbf{#1}, fonttitle=\small,
  breakable, boxrule=0.5pt, left=4pt, right=4pt, top=4pt, bottom=4pt}
\newtcolorbox{compbox}[1][Complexity]{
  colback=compyellow, colframe=compborder,
  title=\textbf{#1}, fonttitle=\small,
  breakable, boxrule=0.5pt, left=4pt, right=4pt, top=4pt, bottom=4pt}
\newtcolorbox{tipbox}[1][Tip]{
  colback=tipgreen, colframe=tipborder,
  title=\textbf{#1}, fonttitle=\small,
  breakable, boxrule=0.5pt, left=4pt, right=4pt, top=4pt, bottom=4pt}
\newtcolorbox{warnbox}[1][Warning]{
  colback=warnred, colframe=warnborder,
  title=\textbf{#1}, fonttitle=\small,
  breakable, boxrule=0.5pt, left=4pt, right=4pt, top=4pt, bottom=4pt}

% ── Page style ───────────────────────────────────────────────
\pagestyle{fancy}
\fancyhf{}
\fancyhead[L]{\small\color{gray}\textit{\leftmark}}
\fancyhead[R]{\small\thepage}
\fancyfoot[C]{\small\color{gray}\textit{DSA in Python --- Complete Reference Guide}}
\renewcommand{\headrulewidth}{0.4pt}
\renewcommand{\footrulewidth}{0.4pt}

% ── Chapter/Section Formatting ───────────────────────────────
\titleformat{\chapter}[display]
  {\normalfont\huge\bfseries\color{accentblue}}
  {\large\chaptertitlename\ \thechapter}{12pt}{\Huge}
\titlespacing*{\chapter}{0pt}{10pt}{20pt}
\titleformat{\section}{\large\bfseries\color{accentblue!80!black}}{{\thesection}}{0.6em}{}
\titleformat{\subsection}{\normalsize\bfseries\color{accentblue!60!black}}{{\thesubsection}}{0.5em}{}

% ── Macros ───────────────────────────────────────────────────
\newcommand{\bigO}[1]{\ensuremath{\mathcal{O}(#1)}}
\newcommand{\py}[1]{\texttt{\color{pyblue}#1}}
\newcommand{\TC}[1]{\textbf{Time:}\ $\mathcal{O}(#1)$\quad}
\newcommand{\SC}[1]{\textbf{Space:}\ $\mathcal{O}(#1)$}

%%%%%%%%%%%%%%%%%%%%%%%%%%%%%%%%%%%%%%%%%%%%%%%%%%%%%%%%%%%%%%%
\begin{document}
%%%%%%%%%%%%%%%%%%%%%%%%%%%%%%%%%%%%%%%%%%%%%%%%%%%%%%%%%%%%%%%

%%── Title Page ──────────────────────────────────────────────
\begin{titlepage}
\centering\vspace*{1.8cm}
{\Huge\bfseries\color{accentblue}Data Structures \&\\[6pt]Algorithms in Python\par}
\vspace{0.5cm}
{\color{gray}\rule{0.65\linewidth}{1pt}\par}
\vspace{0.4cm}
{\Large A Complete Reference for C\texttt{++} Programmers\par}
\vspace{0.3cm}
{\small\color{gray}
Arrays $\cdot$ Linked Lists $\cdot$ Stacks $\cdot$ Queues $\cdot$ Trees $\cdot$ Heaps $\cdot$
Hash Tables $\cdot$ Tries\\
Segment Trees $\cdot$ Fenwick Trees $\cdot$ Sparse Tables $\cdot$ DSU\\
Graphs $\cdot$ Sorting $\cdot$ Searching $\cdot$ Backtracking $\cdot$ Dynamic Programming\\
Greedy $\cdot$ Divide \& Conquer $\cdot$ String Algorithms $\cdot$
Mathematical Algorithms $\cdot$ Bit Manipulation\par}
\vspace{0.5cm}
{\color{gray}\rule{0.65\linewidth}{1pt}\par}
\vspace{1.5cm}
{\large\textit{Every Algorithm $\cdot$ Every Data Structure $\cdot$ Every Complexity}\par}
\vfill{\today}
\end{titlepage}

\tableofcontents\clearpage

%%%%%%%%%%%%%%%%%%%%%%%%%%%%%%%%%%%%%%%%%%%%%%%%%%%%%%%%%%%%%%%
\part{Foundations}
%%%%%%%%%%%%%%%%%%%%%%%%%%%%%%%%%%%%%%%%%%%%%%%%%%%%%%%%%%%%%%%

%%%==========================================================
\chapter{Python for C\texttt{++} Programmers}
%%%==========================================================

This guide assumes full fluency in C\texttt{++} DSA\@. Python is dynamically typed,
garbage-collected, and has arbitrary-precision integers --- no overflow, no manual memory.
The trade-off is speed: expect Python to be 10--100$\times$ slower than C\texttt{++}.

%─────────────────────────────────────────────────────────────
\section{C\texttt{++} $\leftrightarrow$ Python Quick-Reference}
%─────────────────────────────────────────────────────────────

\begin{longtable}{@{}lll@{}}
\toprule
\textbf{Concept} & \textbf{C++} & \textbf{Python 3} \\
\midrule
\endhead
Dynamic array       & \texttt{vector<int>}             & \texttt{list} \\
Stack               & \texttt{stack<T>}                & \texttt{list} or \texttt{collections.deque} \\
Queue               & \texttt{queue<T>}                & \texttt{collections.deque} \\
Deque               & \texttt{deque<T>}                & \texttt{collections.deque} \\
Priority queue (max)& \texttt{priority\_queue<T>}      & \texttt{heapq} (negate for max-heap) \\
Ordered map         & \texttt{map<K,V>}                & \texttt{dict} (insertion-ordered 3.7+) \\
Unordered map       & \texttt{unordered\_map<K,V>}     & \texttt{dict} \\
Ordered set         & \texttt{set<T>}                  & \texttt{sortedcontainers.SortedList} \\
Unordered set       & \texttt{unordered\_set<T>}       & \texttt{set} \\
Pair                & \texttt{pair<A,B>}               & \texttt{tuple} \\
Tuple               & \texttt{tuple<...>}              & \texttt{tuple} \\
String              & \texttt{std::string} (mutable)   & \texttt{str} (immutable) \\
Null pointer        & \texttt{nullptr}                 & \texttt{None} \\
Boolean literals    & \texttt{true / false}            & \texttt{True / False} \\
Integer overflow    & Yes --- use \texttt{long long}   & Never (arbitrary precision) \\
Infinity            & \texttt{INT\_MAX, LLONG\_MAX}    & \texttt{float('inf')} \\
Auto type           & \texttt{auto x = ...}            & No declaration needed \\
Lambda              & \texttt{[](int x)\{return x*2;\}}& \texttt{lambda x: x*2} \\
Range-based for     & \texttt{for(auto x : vec)}       & \texttt{for x in lst} \\
Ternary             & \texttt{cond ? a : b}            & \texttt{a if cond else b} \\
Sort (custom cmp)   & \texttt{sort(a,b,cmp)}           & \texttt{sorted(lst, key=func)} \\
Modular inverse     & Manual (ext. Euclidean)          & \texttt{pow(a, -1, mod)} \\
GCD                 & \texttt{\_\_gcd(a,b)}            & \texttt{math.gcd(a,b)} \\
Swap                & \texttt{swap(a,b)}               & \texttt{a, b = b, a} \\
Bit operations      & \texttt{\&, |, \^{}, $\sim$, <<, >>} & same operators \\
Input               & \texttt{cin >> x}                & \texttt{x = int(input())} \\
\bottomrule
\end{longtable}

%─────────────────────────────────────────────────────────────
\section{Essential Python Idioms}
%─────────────────────────────────────────────────────────────

\begin{lstlisting}
# ── Swap (no temp variable) ───────────────────────────────────
a, b = b, a

# ── Multiple assignment ───────────────────────────────────────
x, y, z = 1, 2, 3
first, *rest = [1, 2, 3, 4]    # first=1, rest=[2,3,4]
*init, last  = [1, 2, 3, 4]    # init=[1,2,3], last=4

# ── List comprehension (replaces C++ for-loop building vector) 
squares = [x**2 for x in range(10)]
evens   = [x for x in range(20) if x % 2 == 0]
flat    = [item for sub in matrix for item in sub]  # flatten 2D

# ── Dict / Set comprehensions ────────────────────────────────
freq  = {ch: s.count(ch) for ch in set(s)}
uniq  = {x for x in arr}

# ── Chained comparisons ───────────────────────────────────────
if 0 <= i < n:   pass       # C++: i >= 0 && i < n

# ── Enumerate (index + value) ─────────────────────────────────
for i, val in enumerate(arr):
    pass

# ── Zip (iterate two sequences in parallel) ──────────────────
for a, b in zip(list1, list2):
    pass
pairs = list(zip(keys, values))

# ── Negative indexing ─────────────────────────────────────────
arr[-1]       # last element
arr[-2]       # second to last

# ── Slicing arr[start:stop:step] ──────────────────────────────
arr[1:4]      # elements at index 1, 2, 3
arr[::-1]     # reversed copy
arr[::2]      # every other element
arr[i:j] = [] # delete elements i to j-1 in-place

# ── String building (use join, NOT +=) ────────────────────────
# SLOW: result = ""; for ch in s: result += ch   # O(n^2)
# FAST:
parts = []
for ch in s:
    parts.append(ch)
result = "".join(parts)    # O(n)

# ── Ternary ───────────────────────────────────────────────────
result = x if condition else y

# ── Truthiness ────────────────────────────────────────────────
# Falsy: None, 0, 0.0, "", [], {}, set()
# All others: truthy

# ── Unpacking in function args ────────────────────────────────
def f(a, b, c): pass
args = [1, 2, 3]
f(*args)              # positional unpacking
kwargs = {'a':1, 'b':2, 'c':3}
f(**kwargs)           # keyword unpacking

# ── walrus operator (:=) assigns and evaluates ────────────────
if (n := len(arr)) > 10:
    print(f"Array has {n} elements")
\end{lstlisting}

%─────────────────────────────────────────────────────────────
\section{Standard Library for DSA}
%─────────────────────────────────────────────────────────────

\begin{lstlisting}
import sys, math, heapq, bisect
from collections import defaultdict, Counter, deque, OrderedDict
from functools import lru_cache, cache, reduce
from itertools import permutations, combinations, combinations_with_replacement, \
                      product, chain, accumulate

# ── sys ──────────────────────────────────────────────────────
sys.setrecursionlimit(300000)      # essential for deep DFS
INF = float('inf')
NINF = float('-inf')
input = sys.stdin.readline         # fast input for competitive programming

# ── math ─────────────────────────────────────────────────────
math.gcd(12, 8)          # 4
math.lcm(4, 6)           # 12  (Python 3.9+)
math.isqrt(17)           # 4  (integer sqrt, no float rounding)
math.log2(1024)          # 10.0
math.comb(10, 3)         # 120  (nCr, exact integer)
math.factorial(10)       # 3628800
math.ceil(3.2)           # 4
math.floor(3.9)          # 3
pow(2, 10, 1000)         # modular exponentiation: 2^10 mod 1000 = 24
pow(3, -1, 7)            # modular inverse: 3^{-1} mod 7 = 5

# ── heapq (min-heap) ─────────────────────────────────────────
h = []
heapq.heappush(h, 5)
top = heapq.heappop(h)           # smallest element
heapq.heapify(lst)               # O(n) in-place
heapq.nlargest(k, lst)           # k largest in O(n log k)
heapq.nsmallest(k, lst)          # k smallest
# Max-heap trick: push negative values
heapq.heappush(h, -val)
max_val = -heapq.heappop(h)

# ── bisect (sorted list binary search) ───────────────────────
arr = [1, 3, 5, 7, 9]
bisect.bisect_left(arr, 5)       # 2  (leftmost insert position)
bisect.bisect_right(arr, 5)      # 3  (rightmost insert position)
bisect.insort(arr, 6)            # insert 6 maintaining sort: O(n)

# ── defaultdict ───────────────────────────────────────────────
graph = defaultdict(list)        # adjacency list
freq  = defaultdict(int)         # frequency map (default 0)
dp    = defaultdict(lambda: INF) # custom default

# ── Counter ───────────────────────────────────────────────────
c = Counter("banana")            # {'a':3,'n':2,'b':1}
c.most_common(2)                 # [('a',3),('n',2)]
c['z']                           # 0 (missing key safe)
c1 + c2                          # union of counters
c1 - c2                          # difference

# ── deque ────────────────────────────────────────────────────
dq = deque([1, 2, 3])
dq.appendleft(0)                 # O(1)  -- front push
dq.popleft()                     # O(1)  -- front pop
dq.rotate(k)                     # rotate right by k
dq.maxlen                        # fixed-size sliding window: deque(maxlen=k)

# ── itertools ────────────────────────────────────────────────
list(permutations([1,2,3]))      # all 3! orderings
list(combinations([1,2,3], 2))   # [(1,2),(1,3),(2,3)]
list(product([0,1], repeat=3))   # all 2^3 binary strings
list(accumulate([1,2,3,4]))      # [1, 3, 6, 10]  (prefix sums)
list(chain([1,2], [3,4]))        # [1, 2, 3, 4]

# ── lru_cache for memoization ────────────────────────────────
@lru_cache(maxsize=None)
def fib(n):
    if n <= 1: return n
    return fib(n-1) + fib(n-2)

# Python 3.9+ shorthand
@cache
def fib(n):
    if n <= 1: return n
    return fib(n-1) + fib(n-2)
\end{lstlisting}

\begin{warnbox}[Recursion Limit]
Python's default recursion limit is 1000. For DFS on large graphs or deep DP recursion, always add \py{sys.setrecursionlimit(300000)} at the top of your file. Alternatively, convert recursive DFS to iterative with an explicit stack.
\end{warnbox}

%%%==========================================================
\chapter{Complexity Analysis}
%%%==========================================================

\section{Big-O Hierarchy}

\[
O(1) < O(\log n) < O(\sqrt{n}) < O(n) < O(n \log n) < O(n^2) < O(n^3) < O(2^n) < O(n!)
\]

Practical rule of thumb for competitive programming: at $10^8$ operations per second,
$n \le 10^8$ tolerates $O(n)$; $n \le 10^4$ tolerates $O(n^2)$; $n \le 500$ tolerates $O(n^3)$;
$n \le 20$ tolerates $O(2^n)$; $n \le 12$ tolerates $O(n!)$.

\section{Python Built-in Complexity Reference}

\begin{longtable}{@{}llll@{}}
\toprule
\textbf{Operation} & \textbf{Container} & \textbf{Avg} & \textbf{Worst} \\
\midrule
\endhead
\texttt{lst[i]}              & list    & $O(1)$        & $O(1)$ \\
\texttt{lst.append(x)}       & list    & $O(1)$ amort  & $O(n)$ \\
\texttt{lst.pop()}           & list    & $O(1)$        & $O(1)$ \\
\texttt{lst.pop(0)}          & list    & $O(n)$        & $O(n)$ \\
\texttt{lst.insert(i,x)}     & list    & $O(n)$        & $O(n)$ \\
\texttt{x in lst}            & list    & $O(n)$        & $O(n)$ \\
\texttt{lst.sort()}          & list    & $O(n\log n)$  & $O(n\log n)$ \\
\texttt{lst[a:b]}            & list    & $O(b{-}a)$    & $O(n)$ \\
\texttt{d[k]}                & dict    & $O(1)$        & $O(n)$ \\
\texttt{d[k] = v}            & dict    & $O(1)$        & $O(n)$ \\
\texttt{k in d}              & dict    & $O(1)$        & $O(n)$ \\
\texttt{del d[k]}            & dict    & $O(1)$        & $O(n)$ \\
\texttt{x in s}              & set     & $O(1)$        & $O(n)$ \\
\texttt{s.add(x)}            & set     & $O(1)$        & $O(n)$ \\
\texttt{s.remove(x)}         & set     & $O(1)$        & $O(n)$ \\
\texttt{s1 | s2} (union)     & set     & $O(m{+}n)$    & --- \\
\texttt{s1 \& s2} (inter.)   & set     & $O(\min(m,n))$& --- \\
\texttt{dq.append()}         & deque   & $O(1)$        & $O(1)$ \\
\texttt{dq.appendleft()}     & deque   & $O(1)$        & $O(1)$ \\
\texttt{dq.pop()}            & deque   & $O(1)$        & $O(1)$ \\
\texttt{dq.popleft()}        & deque   & $O(1)$        & $O(1)$ \\
\texttt{heapq.heappush()}    & heap    & $O(\log n)$   & $O(\log n)$ \\
\texttt{heapq.heappop()}     & heap    & $O(\log n)$   & $O(\log n)$ \\
\texttt{heapq.heapify()}     & heap    & $O(n)$        & $O(n)$ \\
\texttt{bisect.bisect()}     & sorted list & $O(\log n)$ & $O(\log n)$ \\
\texttt{str + str}           & str     & $O(n)$        & $O(n)$ \\
\texttt{"".join(lst)}        & str     & $O(n)$        & $O(n)$ \\
\texttt{len(container)}      & any     & $O(1)$        & $O(1)$ \\
\bottomrule
\end{longtable}

%%%%%%%%%%%%%%%%%%%%%%%%%%%%%%%%%%%%%%%%%%%%%%%%%%%%%%%%%%%%%%%
\part{Linear Data Structures}
%%%%%%%%%%%%%%%%%%%%%%%%%%%%%%%%%%%%%%%%%%%%%%%%%%%%%%%%%%%%%%%

%%%==========================================================
\chapter{Arrays and Strings}
%%%==========================================================

\section{Python Lists --- Dynamic Arrays}

\begin{lstlisting}
# Declaration
arr = []                           # empty, like vector<int>
arr = [0] * n                      # n zeros
arr = list(range(n))               # [0, 1, ..., n-1]
matrix = [[0]*cols for _ in range(rows)]  # 2D --- NOT [[0]*cols]*rows !

# Core operations
arr.append(x)          # push_back  O(1) amortised
arr.pop()              # pop_back   O(1)
arr.pop(i)             # erase(it)  O(n)
arr.insert(i, x)       # insert at i O(n)
arr.remove(x)          # erase first x O(n)
arr.index(x)           # first index of x O(n); raises ValueError if absent
x in arr               # O(n) search; use set for O(1)
arr.sort()             # in-place Timsort O(n log n)
arr.sort(reverse=True)
arr.sort(key=lambda x: x[1])      # sort by second element
sorted(arr)            # returns NEW sorted list (original unchanged)
arr.reverse()          # in-place O(n)
arr.count(x)           # count occurrences O(n)
arr.copy()             # shallow copy O(n)
arr[:]                 # also shallow copy
arr.clear()            # O(n)
arr.extend(other)      # append all elements of other O(k)
\end{lstlisting}

\begin{warnbox}[2D Array Pitfall]
\texttt{[[0]*cols]*rows} creates \emph{rows} references to the \emph{same} inner list --- modifying one row modifies all. Always use \texttt{[[0]*cols for \_ in range(rows)]}.
\end{warnbox}

\section{Strings}

\begin{lstlisting}
s = "hello world"

# Strings are IMMUTABLE -- to mutate, convert to list first
lst = list(s)
lst[0] = 'H'
s = "".join(lst)

# Common methods
s.upper()             # "HELLO WORLD"
s.lower()             # "hello world"
s.strip()             # strip leading/trailing whitespace
s.lstrip(); s.rstrip()
s.split()             # split by whitespace -> list
s.split(',')          # split by delimiter
",".join(["a","b"])   # "a,b"
s.replace('l','r')    # replace ALL occurrences
s.find('lo')          # index of first match, -1 if missing
s.rfind('lo')         # last match
s.count('l')          # count non-overlapping occurrences
s.startswith('he')    # bool
s.endswith('ld')      # bool
s.isalpha()           # all letters?
s.isdigit()           # all digits?
s.isalnum()           # all alphanumeric?
s.isspace()           # all whitespace?
s.zfill(5)            # zero-pad: "007"
ord('a')              # 97   (like (int)'a' in C++)
chr(97)               # 'a'  (like (char)97  in C++)
s[i:j]                # substring, O(j-i)
s[::-1]               # reverse string
\end{lstlisting}

\section{Two Pointers}

\begin{lstlisting}
# ── Two Sum in Sorted Array ── TC: O(n)  SC: O(1) ─────────────
def two_sum_sorted(arr, target):
    lo, hi = 0, len(arr) - 1
    while lo < hi:
        s = arr[lo] + arr[hi]
        if s == target:   return [lo, hi]
        elif s < target:  lo += 1
        else:             hi -= 1
    return []

# ── Three Sum (unique triplets summing to 0) ── O(n^2) ────────
def three_sum(nums):
    nums.sort()
    result = []
    for i in range(len(nums) - 2):
        if i > 0 and nums[i] == nums[i-1]:
            continue
        lo, hi = i + 1, len(nums) - 1
        while lo < hi:
            s = nums[i] + nums[lo] + nums[hi]
            if s == 0:
                result.append([nums[i], nums[lo], nums[hi]])
                while lo < hi and nums[lo] == nums[lo+1]: lo += 1
                while lo < hi and nums[hi] == nums[hi-1]: hi -= 1
                lo += 1; hi -= 1
            elif s < 0: lo += 1
            else:       hi -= 1
    return result

# ── Container with Most Water ── O(n) ────────────────────────
def max_water(height):
    lo, hi = 0, len(height) - 1
    res = 0
    while lo < hi:
        res = max(res, min(height[lo], height[hi]) * (hi - lo))
        if height[lo] < height[hi]: lo += 1
        else:                        hi -= 1
    return res

# ── Trapping Rain Water ── O(n) ────────────────────────────────
def trap(height):
    n = len(height)
    if n == 0: return 0
    left_max  = [0] * n
    right_max = [0] * n
    left_max[0] = height[0]
    for i in range(1, n):
        left_max[i] = max(left_max[i-1], height[i])
    right_max[n-1] = height[n-1]
    for i in range(n-2, -1, -1):
        right_max[i] = max(right_max[i+1], height[i])
    water = 0
    for i in range(n):
        water += min(left_max[i], right_max[i]) - height[i]
    return water

# O(1) space two-pointer version
def trap_two_ptr(height):
    lo, hi = 0, len(height) - 1
    left_max = right_max = 0
    water = 0
    while lo <= hi:
        if height[lo] <= height[hi]:
            if height[lo] >= left_max: left_max = height[lo]
            else:                       water += left_max - height[lo]
            lo += 1
        else:
            if height[hi] >= right_max: right_max = height[hi]
            else:                        water += right_max - height[hi]
            hi -= 1
    return water
\end{lstlisting}

\section{Sliding Window}

\begin{lstlisting}
# ── Fixed Window: Max Sum of Subarray Size k ── O(n) ──────────
def max_sum_k(arr, k):
    window = sum(arr[:k])
    best = window
    for i in range(k, len(arr)):
        window += arr[i] - arr[i-k]
        best = max(best, window)
    return best

# ── Variable Window: Longest Substring Without Repeating ──────
def longest_unique_substr(s):        # TC: O(n)  SC: O(1)
    seen = {}
    lo = best = 0
    for hi, ch in enumerate(s):
        if ch in seen and seen[ch] >= lo:
            lo = seen[ch] + 1
        seen[ch] = hi
        best = max(best, hi - lo + 1)
    return best

# ── Minimum Window Substring ── O(|s|+|t|) ────────────────────
from collections import Counter
def min_window(s, t):
    need = Counter(t)
    have = total = 0
    total = len(need)
    window = {}
    lo = 0
    res_len = float('inf')
    res = ""
    for hi, ch in enumerate(s):
        window[ch] = window.get(ch, 0) + 1
        if ch in need and window[ch] == need[ch]:
            have += 1
        while have == total:
            if hi - lo + 1 < res_len:
                res_len = hi - lo + 1
                res = s[lo:hi+1]
            window[s[lo]] -= 1
            if s[lo] in need and window[s[lo]] < need[s[lo]]:
                have -= 1
            lo += 1
    return res

# ── Sliding Window Maximum (deque) ── O(n) ────────────────────
from collections import deque
def max_sliding_window(nums, k):
    dq = deque()       # stores indices, decreasing values
    result = []
    for i, x in enumerate(nums):
        # remove indices out of window
        while dq and dq[0] < i - k + 1:
            dq.popleft()
        # remove smaller elements from back
        while dq and nums[dq[-1]] < x:
            dq.pop()
        dq.append(i)
        if i >= k - 1:
            result.append(nums[dq[0]])
    return result
\end{lstlisting}

\section{Prefix Sums}

\begin{lstlisting}
# ── 1D Prefix Sum ─────────────────────────────────────────────
def build_prefix(arr):              # TC: O(n)  SC: O(n)
    n = len(arr)
    prefix = [0] * (n + 1)
    for i in range(n):
        prefix[i+1] = prefix[i] + arr[i]
    return prefix
# Range sum [l, r] in O(1)
def range_sum(prefix, l, r): return prefix[r+1] - prefix[l]

# Using itertools.accumulate (Pythonic)
from itertools import accumulate
prefix = list(accumulate(arr, initial=0))   # [0, a0, a0+a1, ...]

# ── Subarray Sum Equals k ── O(n) ─────────────────────────────
def subarray_sum_k(nums, k):
    count = curr = 0
    freq = {0: 1}
    for x in nums:
        curr += x
        count += freq.get(curr - k, 0)
        freq[curr] = freq.get(curr, 0) + 1
    return count

# ── Product of Array Except Self ── O(n)  SC: O(1) output excl
def product_except_self(nums):
    n = len(nums)
    res = [1] * n
    left = 1
    for i in range(n):
        res[i] = left; left *= nums[i]
    right = 1
    for i in range(n-1, -1, -1):
        res[i] *= right; right *= nums[i]
    return res

# ── 2D Prefix Sum ─────────────────────────────────────────────
def build_2d_prefix(mat):
    m, n = len(mat), len(mat[0])
    P = [[0]*(n+1) for _ in range(m+1)]
    for i in range(1, m+1):
        for j in range(1, n+1):
            P[i][j] = mat[i-1][j-1] + P[i-1][j] + P[i][j-1] - P[i-1][j-1]
    return P

def rect_sum(P, r1, c1, r2, c2):  # O(1) query
    return P[r2+1][c2+1] - P[r1][c2+1] - P[r2+1][c1] + P[r1][c1]
\end{lstlisting}

\section{Kadane's Algorithm}

\begin{lstlisting}
# ── Maximum Subarray Sum ── O(n)  SC: O(1) ────────────────────
def max_subarray(nums):
    best = curr = nums[0]
    for x in nums[1:]:
        curr = max(x, curr + x)
        best = max(best, curr)
    return best

# ── With Indices ──────────────────────────────────────────────
def max_subarray_indices(nums):
    best = curr = nums[0]
    start = end = temp_start = 0
    for i in range(1, len(nums)):
        if nums[i] > curr + nums[i]:
            curr = nums[i]; temp_start = i
        else:
            curr += nums[i]
        if curr > best:
            best = curr; start = temp_start; end = i
    return best, start, end

# ── Maximum Circular Subarray Sum ── O(n) ────────────────────
def max_circular(nums):
    def kadane_max(a):
        cur = res = a[0]
        for x in a[1:]: cur = max(x, cur+x); res = max(res, cur)
        return res
    def kadane_min(a):
        cur = res = a[0]
        for x in a[1:]: cur = min(x, cur+x); res = min(res, cur)
        return res
    max_straight = kadane_max(nums)
    circular     = sum(nums) - kadane_min(nums)
    return max(max_straight, circular) if circular != 0 else max_straight
\end{lstlisting}

\section{Other Important Array Algorithms}

\begin{lstlisting}
# ── Boyer-Moore Majority Vote ── O(n)  SC: O(1) ───────────────
def majority_element(nums):     # element appearing > n/2 times
    candidate = count = 0
    for x in nums:
        if count == 0: candidate = x
        count += 1 if x == candidate else -1
    return candidate

# ── Dutch National Flag (3-way partition) ── O(n) ─────────────
def sort_colors(nums):          # sort 0s, 1s, 2s in one pass
    lo = mid = 0; hi = len(nums) - 1
    while mid <= hi:
        if   nums[mid] == 0: nums[lo], nums[mid] = nums[mid], nums[lo]; lo += 1; mid += 1
        elif nums[mid] == 1: mid += 1
        else:                 nums[mid], nums[hi] = nums[hi], nums[mid]; hi -= 1

# ── Rotate Array Right by k ── O(n)  SC: O(1) ────────────────
def rotate(nums, k):
    n = len(nums); k %= n
    def rev(l, r):
        while l < r: nums[l], nums[r] = nums[r], nums[l]; l += 1; r -= 1
    rev(0, n-1); rev(0, k-1); rev(k, n-1)

# ── Next Permutation ── O(n) ──────────────────────────────────
def next_permutation(nums):
    n = len(nums); i = n - 2
    while i >= 0 and nums[i] >= nums[i+1]: i -= 1
    if i >= 0:
        j = n - 1
        while nums[j] <= nums[i]: j -= 1
        nums[i], nums[j] = nums[j], nums[i]
    lo, hi = i+1, n-1
    while lo < hi: nums[lo], nums[hi] = nums[hi], nums[lo]; lo += 1; hi -= 1

# ── Find All Duplicates (using sign trick) ── O(n)  SC: O(1) ──
def find_duplicates(nums):      # elements in [1, n]
    result = []
    for x in nums:
        idx = abs(x) - 1
        if nums[idx] < 0: result.append(abs(x))
        else:              nums[idx] = -nums[idx]
    return result

# ── Longest Consecutive Sequence ── O(n) ─────────────────────
def longest_consecutive(nums):
    num_set = set(nums)
    best = 0
    for x in num_set:
        if x - 1 not in num_set:   # start of a sequence
            cur = x; length = 1
            while cur + 1 in num_set: cur += 1; length += 1
            best = max(best, length)
    return best
\end{lstlisting}


%%%==========================================================
\chapter{Linked Lists}
%%%==========================================================

\section{Node Class and Singly Linked List}

\begin{lstlisting}
class ListNode:
    def __init__(self, val=0, next=None):
        self.val  = val
        self.next = next

class SinglyLinkedList:
    def __init__(self):
        self.head = None
        self.size = 0

    # ── Insert at head ── O(1) ──────────────────────────────
    def push_front(self, val):
        self.head = ListNode(val, self.head)
        self.size += 1

    # ── Insert at tail ── O(n) or O(1) with tail pointer ───
    def push_back(self, val):
        if not self.head:
            self.head = ListNode(val); self.size += 1; return
        cur = self.head
        while cur.next: cur = cur.next
        cur.next = ListNode(val)
        self.size += 1

    # ── Insert at position k (0-indexed) ── O(k) ───────────
    def insert(self, k, val):
        if k == 0: self.push_front(val); return
        cur = self.head
        for _ in range(k - 1):
            if not cur: return
            cur = cur.next
        node = ListNode(val, cur.next)
        cur.next = node
        self.size += 1

    # ── Delete first node with given value ── O(n) ─────────
    def delete(self, val):
        dummy = ListNode(0, self.head)
        cur = dummy
        while cur.next:
            if cur.next.val == val:
                cur.next = cur.next.next; self.size -= 1; break
            cur = cur.next
        self.head = dummy.next

    # ── Search ── O(n) ──────────────────────────────────────
    def search(self, val):
        cur = self.head
        while cur:
            if cur.val == val: return True
            cur = cur.next
        return False

    # ── Reverse in-place ── O(n)  SC: O(1) ─────────────────
    def reverse(self):
        prev = None; cur = self.head
        while cur:
            nxt = cur.next
            cur.next = prev
            prev = cur; cur = nxt
        self.head = prev

    # ── Print ───────────────────────────────────────────────
    def to_list(self):
        res = []; cur = self.head
        while cur: res.append(cur.val); cur = cur.next
        return res
\end{lstlisting}

\section{Doubly Linked List}

\begin{lstlisting}
class DNode:
    def __init__(self, val=0, prev=None, next=None):
        self.val  = val
        self.prev = prev
        self.next = next

class DoublyLinkedList:
    def __init__(self):
        # sentinel head and tail (dummy nodes)
        self.head = DNode()
        self.tail = DNode()
        self.head.next = self.tail
        self.tail.prev = self.head
        self.size = 0

    def _insert_after(self, node, val):   # internal helper
        new = DNode(val, node, node.next)
        node.next.prev = new
        node.next = new
        self.size += 1

    def push_front(self, val): self._insert_after(self.head, val)
    def push_back(self, val):  self._insert_after(self.tail.prev, val)

    def _remove(self, node):              # internal helper
        node.prev.next = node.next
        node.next.prev = node.prev
        self.size -= 1

    def pop_front(self):
        if self.head.next == self.tail: return None
        val = self.head.next.val
        self._remove(self.head.next); return val

    def pop_back(self):
        if self.tail.prev == self.head: return None
        val = self.tail.prev.val
        self._remove(self.tail.prev); return val
\end{lstlisting}

\section{Common Linked List Algorithms}

\begin{lstlisting}
# ── Reverse a Linked List ── O(n)  SC: O(1) ───────────────────
def reverse_list(head):
    prev = None; cur = head
    while cur:
        nxt = cur.next; cur.next = prev; prev = cur; cur = nxt
    return prev

# ── Find Middle (slow-fast pointers) ── O(n)  SC: O(1) ────────
def find_middle(head):
    slow = fast = head
    while fast and fast.next:
        slow = slow.next; fast = fast.next.next
    return slow      # for even length: returns second middle

# ── Floyd's Cycle Detection ── O(n)  SC: O(1) ─────────────────
def has_cycle(head):
    slow = fast = head
    while fast and fast.next:
        slow = slow.next; fast = fast.next.next
        if slow is fast: return True
    return False

def find_cycle_start(head):     # start of cycle or None
    slow = fast = head
    while fast and fast.next:
        slow = slow.next; fast = fast.next.next
        if slow is fast:
            # move one pointer to head; advance both at speed 1
            slow = head
            while slow is not fast:
                slow = slow.next; fast = fast.next
            return slow
    return None

# ── Merge Two Sorted Lists ── O(m+n)  SC: O(1) ────────────────
def merge_two_sorted(l1, l2):
    dummy = cur = ListNode(0)
    while l1 and l2:
        if l1.val <= l2.val: cur.next = l1; l1 = l1.next
        else:                 cur.next = l2; l2 = l2.next
        cur = cur.next
    cur.next = l1 if l1 else l2
    return dummy.next

# ── Remove nth Node from End ── O(n)  SC: O(1) ────────────────
def remove_nth_from_end(head, n):
    dummy = ListNode(0, head)
    fast = slow = dummy
    for _ in range(n + 1): fast = fast.next
    while fast:
        fast = fast.next; slow = slow.next
    slow.next = slow.next.next
    return dummy.next

# ── Reverse Nodes in k-Groups ── O(n)  SC: O(k) call stack ────
def reverse_k_group(head, k):
    def reverse(node, k):
        prev = None; cur = node; count = 0
        while cur and count < k:
            nxt = cur.next; cur.next = prev; prev = cur
            cur = nxt; count += 1
        return prev, node, cur   # (new_head, old_head/new_tail, next_group_head)

    def count_nodes(node):
        c = 0
        while node: c += 1; node = node.next
        return c

    if count_nodes(head) < k: return head
    new_head, tail, nxt = reverse(head, k)
    tail.next = reverse_k_group(nxt, k)
    return new_head

# ── Intersection of Two Linked Lists ── O(m+n)  SC: O(1) ──────
def get_intersection(headA, headB):
    a, b = headA, headB
    while a is not b:
        a = a.next if a else headB
        b = b.next if b else headA
    return a    # None if no intersection

# ── Add Two Numbers (digits as linked lists) ── O(max(m,n)) ───
def add_two_numbers(l1, l2):
    dummy = cur = ListNode(0)
    carry = 0
    while l1 or l2 or carry:
        v1 = l1.val if l1 else 0
        v2 = l2.val if l2 else 0
        total = v1 + v2 + carry
        carry = total // 10
        cur.next = ListNode(total % 10)
        cur = cur.next
        if l1: l1 = l1.next
        if l2: l2 = l2.next
    return dummy.next

# ── Flatten a Multilevel Doubly Linked List ─────────────────
def flatten(head):
    cur = head
    while cur:
        if cur.child:
            nxt = cur.next
            cur.next = cur.child
            cur.child.prev = cur
            cur.child = None
            tail = cur.next
            while tail.next: tail = tail.next
            tail.next = nxt
            if nxt: nxt.prev = tail
        cur = cur.next
    return head

# ── Sort Linked List (Merge Sort) ── O(n log n)  SC: O(log n) ─
def sort_list(head):
    if not head or not head.next: return head
    mid = find_middle(head)
    second = mid.next; mid.next = None
    left  = sort_list(head)
    right = sort_list(second)
    return merge_two_sorted(left, right)

# ── LRU Cache using OrderedDict ── all ops O(1) ────────────────
from collections import OrderedDict
class LRUCache:
    def __init__(self, capacity):
        self.cap = capacity
        self.cache = OrderedDict()    # key -> value, MRU at end

    def get(self, key):
        if key not in self.cache: return -1
        self.cache.move_to_end(key)
        return self.cache[key]

    def put(self, key, value):
        if key in self.cache: self.cache.move_to_end(key)
        self.cache[key] = value
        if len(self.cache) > self.cap:
            self.cache.popitem(last=False)    # remove LRU (front)
\end{lstlisting}

%%%==========================================================
\chapter{Stacks}
%%%==========================================================

\section{Stack Implementations}

\begin{lstlisting}
# ── Using list (most common) ──────────────────────────────────
stack = []
stack.append(x)        # push  O(1)
stack.pop()            # pop   O(1)  -- raises IndexError if empty
stack[-1]              # peek  O(1)
len(stack) == 0        # is_empty

# ── Using deque (slightly faster for large data) ──────────────
from collections import deque
stack = deque()
stack.append(x)        # push
stack.pop()            # pop
stack[-1]              # peek

# ── Stack class wrapper ───────────────────────────────────────
class Stack:
    def __init__(self):         self._data = []
    def push(self, x):          self._data.append(x)
    def pop(self):              return self._data.pop()
    def peek(self):             return self._data[-1]
    def is_empty(self):         return len(self._data) == 0
    def __len__(self):          return len(self._data)
\end{lstlisting}

\section{Classic Stack Problems}

\begin{lstlisting}
# ── Valid Parentheses ── O(n)  SC: O(n) ───────────────────────
def is_valid(s):
    match = {')':'(', ']':'[', '}':'{'}
    stack = []
    for ch in s:
        if ch in '([{':
            stack.append(ch)
        else:
            if not stack or stack[-1] != match[ch]:
                return False
            stack.pop()
    return len(stack) == 0

# ── Min Stack (push, pop, peek, getMin all O(1)) ─────────────
class MinStack:
    def __init__(self):
        self.stack = []
        self.min_stack = []

    def push(self, val):
        self.stack.append(val)
        if not self.min_stack or val <= self.min_stack[-1]:
            self.min_stack.append(val)

    def pop(self):
        val = self.stack.pop()
        if val == self.min_stack[-1]:
            self.min_stack.pop()

    def top(self):      return self.stack[-1]
    def getMin(self):   return self.min_stack[-1]

# ── Infix to Postfix (Shunting-Yard) ── O(n) ─────────────────
def infix_to_postfix(expr):
    prec = {'+':1, '-':1, '*':2, '/':2, '^':3}
    assoc = {'^': 'R'}          # all others left-associative
    output = []; stack = []
    for token in expr.split():
        if token.lstrip('-').isdigit():
            output.append(token)
        elif token == '(':
            stack.append(token)
        elif token == ')':
            while stack and stack[-1] != '(':
                output.append(stack.pop())
            stack.pop()
        else:   # operator
            while (stack and stack[-1] != '(' and
                   stack[-1] in prec and
                   (prec[stack[-1]] > prec[token] or
                    (prec[stack[-1]] == prec[token] and assoc.get(token) != 'R'))):
                output.append(stack.pop())
            stack.append(token)
    while stack: output.append(stack.pop())
    return " ".join(output)

# ── Evaluate Postfix ── O(n) ──────────────────────────────────
def eval_postfix(expr):
    stack = []
    for token in expr.split():
        if token.lstrip('-').isdigit():
            stack.append(int(token))
        else:
            b, a = stack.pop(), stack.pop()
            if   token == '+': stack.append(a + b)
            elif token == '-': stack.append(a - b)
            elif token == '*': stack.append(a * b)
            elif token == '/': stack.append(int(a / b))
    return stack[0]
\end{lstlisting}

\section{Monotonic Stack}

A monotonic stack maintains elements in either increasing or decreasing order.
It solves \textbf{Next Greater/Smaller Element}-style problems in $O(n)$.

\begin{lstlisting}
# ── Next Greater Element (to the right) ── O(n) ───────────────
def next_greater(arr):
    n = len(arr)
    result = [-1] * n
    stack = []   # stores indices, values decrease top-to-bottom
    for i, x in enumerate(arr):
        while stack and arr[stack[-1]] < x:
            idx = stack.pop()
            result[idx] = x
        stack.append(i)
    return result

# ── Previous Smaller Element ── O(n) ─────────────────────────
def prev_smaller(arr):
    n = len(arr)
    result = [-1] * n
    stack = []
    for i, x in enumerate(arr):
        while stack and arr[stack[-1]] >= x:
            stack.pop()
        result[i] = arr[stack[-1]] if stack else -1
        stack.append(i)
    return result

# ── Stock Span Problem ── O(n) ────────────────────────────────
def stock_span(prices):
    stack = []   # stores (price, span)
    result = []
    for p in prices:
        span = 1
        while stack and stack[-1][0] <= p:
            span += stack.pop()[1]
        stack.append((p, span))
        result.append(span)
    return result

# ── Largest Rectangle in Histogram ── O(n)  SC: O(n) ─────────
def largest_rectangle(heights):
    heights = heights + [0]   # sentinel
    stack = [-1]              # stack of indices
    max_area = 0
    for i, h in enumerate(heights):
        while stack[-1] != -1 and heights[stack[-1]] >= h:
            height = heights[stack.pop()]
            width  = i - stack[-1] - 1
            max_area = max(max_area, height * width)
        stack.append(i)
    return max_area

# ── Maximal Rectangle (in binary matrix) ── O(m*n) ───────────
def maximal_rectangle(matrix):
    if not matrix: return 0
    n = len(matrix[0])
    heights = [0] * n
    max_area = 0
    for row in matrix:
        for j in range(n):
            heights[j] = heights[j] + 1 if row[j] == '1' else 0
        max_area = max(max_area, largest_rectangle(heights))
    return max_area

# ── Daily Temperatures (Next Warmer Day) ── O(n) ──────────────
def daily_temperatures(T):
    n = len(T)
    result = [0] * n
    stack = []
    for i, t in enumerate(T):
        while stack and T[stack[-1]] < t:
            idx = stack.pop()
            result[idx] = i - idx
        stack.append(i)
    return result

# ── Remove k Digits to Get Smallest Number ── O(n) ───────────
def remove_k_digits(num, k):
    stack = []
    for d in num:
        while k and stack and stack[-1] > d:
            stack.pop(); k -= 1
        stack.append(d)
    # Remove remaining k digits from the end
    stack = stack[:len(stack)-k]
    # Remove leading zeros
    result = "".join(stack).lstrip('0')
    return result if result else "0"

# ── Score of Parentheses ── O(n)  SC: O(n) ────────────────────
def score_of_parentheses(s):
    stack = [0]
    for ch in s:
        if ch == '(':
            stack.append(0)
        else:
            v = stack.pop()
            stack[-1] += max(2 * v, 1)
    return stack[0]

# ── Decode String (e.g. "3[a2[c]]" -> "accaccacc") ── O(n) ───
def decode_string(s):
    stack = []
    cur_str = ""
    cur_num = 0
    for ch in s:
        if ch.isdigit():
            cur_num = cur_num * 10 + int(ch)
        elif ch == '[':
            stack.append((cur_str, cur_num))
            cur_str = ""; cur_num = 0
        elif ch == ']':
            prev_str, num = stack.pop()
            cur_str = prev_str + cur_str * num
        else:
            cur_str += ch
    return cur_str
\end{lstlisting}

\begin{compbox}[Stack Complexities]
\begin{tabular}{@{}lll@{}}
\toprule
\textbf{Operation} & \textbf{Time} & \textbf{Space} \\
\midrule
Push / Pop / Peek & $O(1)$ & --- \\
Valid Parentheses & $O(n)$ & $O(n)$ \\
Next Greater Element & $O(n)$ & $O(n)$ \\
Largest Rectangle & $O(n)$ & $O(n)$ \\
\bottomrule
\end{tabular}
\end{compbox}

%%%==========================================================
\chapter{Queues and Deques}
%%%==========================================================

\section{Queue}

\begin{lstlisting}
from collections import deque

# ── Queue using deque ── all ops O(1) ─────────────────────────
queue = deque()
queue.append(x)        # enqueue (rear)
queue.popleft()        # dequeue (front)
queue[0]               # front element (peek)
len(queue) == 0        # is_empty

# ── queue.Queue (thread-safe) ─────────────────────────────────
import queue
q = queue.Queue(maxsize=10)    # bounded queue
q.put(x)                       # blocks if full
q.get()                        # blocks if empty
q.empty()
q.full()

# ── Circular Queue implementation ────────────────────────────
class CircularQueue:
    def __init__(self, k):
        self.k = k
        self.q = [0] * k
        self.head = self.tail = -1
        self.size = 0

    def enqueue(self, val):
        if self.is_full(): return False
        self.tail = (self.tail + 1) % self.k
        self.q[self.tail] = val
        if self.head == -1: self.head = 0
        self.size += 1; return True

    def dequeue(self):
        if self.is_empty(): return False
        if self.head == self.tail: self.head = self.tail = -1
        else: self.head = (self.head + 1) % self.k
        self.size -= 1; return True

    def front(self):    return self.q[self.head] if not self.is_empty() else -1
    def rear(self):     return self.q[self.tail] if not self.is_empty() else -1
    def is_empty(self): return self.size == 0
    def is_full(self):  return self.size == self.k
\end{lstlisting}

\section{Priority Queue}

\begin{lstlisting}
import heapq

# ── Min-heap (default) ────────────────────────────────────────
h = []
heapq.heappush(h, 3)
heapq.heappush(h, 1)
heapq.heappush(h, 2)
print(heapq.heappop(h))    # 1

# ── Max-heap (negate values) ──────────────────────────────────
h = []
for x in [3, 1, 4, 1, 5]:
    heapq.heappush(h, -x)
print(-heapq.heappop(h))   # 5

# ── Heap with (priority, value) tuples ────────────────────────
tasks = []
heapq.heappush(tasks, (2, "low priority"))
heapq.heappush(tasks, (0, "urgent"))
priority, task = heapq.heappop(tasks)   # (0, "urgent")

# ── Build heap from list ── O(n) ──────────────────────────────
arr = [3, 1, 4, 1, 5, 9, 2, 6]
heapq.heapify(arr)           # in-place min-heap

# ── K-th Largest Element ── O(n log k) ───────────────────────
def kth_largest(nums, k):
    h = []
    for x in nums:
        heapq.heappush(h, x)
        if len(h) > k: heapq.heappop(h)
    return h[0]

# ── Merge K Sorted Lists ── O(n log k) ────────────────────────
def merge_k_sorted(lists):
    h = []
    for i, lst in enumerate(lists):
        if lst: heapq.heappush(h, (lst[0], i, 0))
    result = []
    while h:
        val, i, j = heapq.heappop(h)
        result.append(val)
        if j + 1 < len(lists[i]):
            heapq.heappush(h, (lists[i][j+1], i, j+1))
    return result

# ── Median of Data Stream ─────────────────────────────────────
class MedianFinder:
    def __init__(self):
        self.lo = []   # max-heap (negated) for lower half
        self.hi = []   # min-heap for upper half

    def add_num(self, num):
        heapq.heappush(self.lo, -num)
        # Balance: lo's max must <= hi's min
        if self.hi and -self.lo[0] > self.hi[0]:
            heapq.heappush(self.hi, -heapq.heappop(self.lo))
        # Keep sizes equal or lo has one extra
        if len(self.lo) > len(self.hi) + 1:
            heapq.heappush(self.hi, -heapq.heappop(self.lo))
        elif len(self.hi) > len(self.lo):
            heapq.heappush(self.lo, -heapq.heappop(self.hi))

    def find_median(self):
        if len(self.lo) > len(self.hi): return float(-self.lo[0])
        return (-self.lo[0] + self.hi[0]) / 2.0

# ── Task Scheduler ── O(n log n) ──────────────────────────────
from collections import Counter
import heapq
def least_interval(tasks, n):
    freq = Counter(tasks)
    max_heap = [-f for f in freq.values()]
    heapq.heapify(max_heap)
    time = 0
    queue = []   # (available_time, -count)
    while max_heap or queue:
        if queue and queue[0][0] <= time:
            _, count = queue.pop(0)
            heapq.heappush(max_heap, count)
        if max_heap:
            count = heapq.heappop(max_heap) + 1   # +1 since negated
            if count < 0:
                queue.append((time + n + 1, count))
        time += 1
    return time
\end{lstlisting}


%%%%%%%%%%%%%%%%%%%%%%%%%%%%%%%%%%%%%%%%%%%%%%%%%%%%%%%%%%%%%%%
\part{Non-Linear Data Structures}
%%%%%%%%%%%%%%%%%%%%%%%%%%%%%%%%%%%%%%%%%%%%%%%%%%%%%%%%%%%%%%%

%%%==========================================================
\chapter{Trees}
%%%==========================================================

\section{Binary Tree --- Traversals}

\begin{lstlisting}
class TreeNode:
    def __init__(self, val=0, left=None, right=None):
        self.val   = val
        self.left  = left
        self.right = right

# ── Recursive Traversals ─────────────────────────────────────
def inorder(root):      # left, root, right
    if not root: return []
    return inorder(root.left) + [root.val] + inorder(root.right)

def preorder(root):     # root, left, right
    if not root: return []
    return [root.val] + preorder(root.left) + preorder(root.right)

def postorder(root):    # left, right, root
    if not root: return []
    return postorder(root.left) + postorder(root.right) + [root.val]

# ── Iterative Inorder ── O(n)  SC: O(h) ──────────────────────
def inorder_iter(root):
    stack = []; result = []; cur = root
    while cur or stack:
        while cur: stack.append(cur); cur = cur.left
        cur = stack.pop()
        result.append(cur.val)
        cur = cur.right
    return result

# ── Iterative Preorder ── O(n) ────────────────────────────────
def preorder_iter(root):
    if not root: return []
    stack = [root]; result = []
    while stack:
        node = stack.pop()
        result.append(node.val)
        if node.right: stack.append(node.right)
        if node.left:  stack.append(node.left)
    return result

# ── Iterative Postorder ── O(n) ───────────────────────────────
def postorder_iter(root):
    if not root: return []
    stack = [root]; result = []
    while stack:
        node = stack.pop()
        result.append(node.val)
        if node.left:  stack.append(node.left)
        if node.right: stack.append(node.right)
    return result[::-1]

# ── Level Order (BFS) ── O(n)  SC: O(w) w=max width ──────────
from collections import deque
def level_order(root):
    if not root: return []
    queue = deque([root]); result = []
    while queue:
        level = []
        for _ in range(len(queue)):
            node = queue.popleft()
            level.append(node.val)
            if node.left:  queue.append(node.left)
            if node.right: queue.append(node.right)
        result.append(level)
    return result

# ── Morris Inorder (O(1) space, no stack/recursion) ── O(n) ──
def morris_inorder(root):
    result = []; cur = root
    while cur:
        if not cur.left:
            result.append(cur.val); cur = cur.right
        else:
            prev = cur.left
            while prev.right and prev.right is not cur:
                prev = prev.right
            if not prev.right:
                prev.right = cur; cur = cur.left
            else:
                prev.right = None; result.append(cur.val); cur = cur.right
    return result
\end{lstlisting}

\section{Binary Tree Properties}

\begin{lstlisting}
# ── Height ── O(n) ────────────────────────────────────────────
def height(root):
    if not root: return 0
    return 1 + max(height(root.left), height(root.right))

# ── Diameter (longest path between any two nodes) ── O(n) ────
def diameter(root):
    max_d = [0]
    def dfs(node):
        if not node: return 0
        l = dfs(node.left); r = dfs(node.right)
        max_d[0] = max(max_d[0], l + r)
        return 1 + max(l, r)
    dfs(root); return max_d[0]

# ── Check if Balanced ── O(n) ────────────────────────────────
def is_balanced(root):
    def check(node):
        if not node: return 0
        l = check(node.left);  if l == -1: return -1
        r = check(node.right); if r == -1: return -1
        return -1 if abs(l - r) > 1 else 1 + max(l, r)
    return check(root) != -1

# ── Count Nodes ── O(log^2 n) for complete tree ────────────── 
def count_nodes(root):
    if not root: return 0
    lh = rh = 0
    l = r = root
    while l: lh += 1; l = l.left
    while r: rh += 1; r = r.right
    if lh == rh: return (1 << lh) - 1   # perfect tree
    return 1 + count_nodes(root.left) + count_nodes(root.right)

# ── Lowest Common Ancestor ── O(n) ───────────────────────────
def lca(root, p, q):
    if not root or root is p or root is q: return root
    left  = lca(root.left, p, q)
    right = lca(root.right, p, q)
    if left and right: return root
    return left if left else right

# ── Maximum Path Sum ── O(n) ──────────────────────────────────
def max_path_sum(root):
    best = [float('-inf')]
    def dfs(node):
        if not node: return 0
        l = max(0, dfs(node.left))
        r = max(0, dfs(node.right))
        best[0] = max(best[0], node.val + l + r)
        return node.val + max(l, r)
    dfs(root); return best[0]

# ── Serialize / Deserialize ── O(n) ───────────────────────────
def serialize(root):
    if not root: return "null"
    return str(root.val) + "," + serialize(root.left) + "," + serialize(root.right)

def deserialize(data):
    vals = iter(data.split(","))
    def build():
        v = next(vals)
        if v == "null": return None
        node = TreeNode(int(v))
        node.left  = build()
        node.right = build()
        return node
    return build()

# ── Check if Two Trees are Identical ── O(n) ─────────────────
def is_same(p, q):
    if not p and not q: return True
    if not p or not q or p.val != q.val: return False
    return is_same(p.left, q.left) and is_same(p.right, q.right)

# ── Invert (Mirror) a Binary Tree ── O(n) ────────────────────
def invert_tree(root):
    if not root: return None
    root.left, root.right = invert_tree(root.right), invert_tree(root.left)
    return root

# ── Right Side View (rightmost node at each level) ── O(n) ───
def right_side_view(root):
    result = []; queue = deque([root]) if root else deque()
    while queue:
        for i in range(len(queue)):
            node = queue.popleft()
            if i == 0: result.append(node.val)   # last in level (reversed BFS) 
            # OR: rightmost = last seen at each level using BFS left-to-right
        # Re-do: standard right-side-view
    return result

def right_side_view_v2(root):
    if not root: return []
    result = []; queue = deque([root])
    while queue:
        for i in range(len(queue)):
            node = queue.popleft()
            if i == len(queue): result.append(node.val)  # this logic is wrong
            if node.left:  queue.append(node.left)
            if node.right: queue.append(node.right)
    return result

# Correct right side view:
def right_side_view(root):
    if not root: return []
    result = []; queue = deque([root])
    while queue:
        level_size = len(queue)
        for i in range(level_size):
            node = queue.popleft()
            if i == level_size - 1: result.append(node.val)
            if node.left:  queue.append(node.left)
            if node.right: queue.append(node.right)
    return result
\end{lstlisting}

\section{Binary Search Tree (BST)}

\begin{lstlisting}
class BST:
    def __init__(self): self.root = None

    # ── Insert ── O(h) avg O(log n), worst O(n) ──────────────
    def insert(self, val):
        def _insert(node, val):
            if not node: return TreeNode(val)
            if val < node.val:   node.left  = _insert(node.left, val)
            elif val > node.val: node.right = _insert(node.right, val)
            return node
        self.root = _insert(self.root, val)

    # ── Search ── O(h) ────────────────────────────────────────
    def search(self, val):
        cur = self.root
        while cur:
            if   val == cur.val: return cur
            elif val < cur.val:  cur = cur.left
            else:                cur = cur.right
        return None

    # ── Delete ── O(h) ────────────────────────────────────────
    def delete(self, val):
        def _delete(node, val):
            if not node: return None
            if val < node.val:
                node.left = _delete(node.left, val)
            elif val > node.val:
                node.right = _delete(node.right, val)
            else:
                if not node.left:  return node.right
                if not node.right: return node.left
                # Find inorder successor (min of right subtree)
                succ = node.right
                while succ.left: succ = succ.left
                node.val = succ.val
                node.right = _delete(node.right, succ.val)
            return node
        self.root = _delete(self.root, val)

    # ── Floor and Ceil ── O(h) ───────────────────────────────
    def floor(self, val):      # largest value <= val
        res = None; cur = self.root
        while cur:
            if   cur.val == val: return val
            elif cur.val < val:  res = cur.val; cur = cur.right
            else:                cur = cur.left
        return res

    def ceil(self, val):       # smallest value >= val
        res = None; cur = self.root
        while cur:
            if   cur.val == val: return val
            elif cur.val > val:  res = cur.val; cur = cur.left
            else:                cur = cur.right
        return res

# ── Validate BST ── O(n)  SC: O(h) ───────────────────────────
def is_valid_bst(root, lo=float('-inf'), hi=float('inf')):
    if not root: return True
    if not (lo < root.val < hi): return False
    return (is_valid_bst(root.left, lo, root.val) and
            is_valid_bst(root.right, root.val, hi))

# ── Kth Smallest Element ── O(k)  SC: O(h) ───────────────────
def kth_smallest(root, k):
    stack = []; cur = root
    while True:
        while cur: stack.append(cur); cur = cur.left
        cur = stack.pop(); k -= 1
        if k == 0: return cur.val
        cur = cur.right

# ── Inorder Successor ── O(h) ────────────────────────────────
def inorder_successor(root, p):
    succ = None
    while root:
        if p.val < root.val: succ = root; root = root.left
        else:                              root = root.right
    return succ

# ── Convert Sorted Array to Balanced BST ── O(n) ─────────────
def sorted_array_to_bst(nums):
    if not nums: return None
    mid  = len(nums) // 2
    node = TreeNode(nums[mid])
    node.left  = sorted_array_to_bst(nums[:mid])
    node.right = sorted_array_to_bst(nums[mid+1:])
    return node

# ── BST to Sorted Doubly Linked List ── O(n) ─────────────────
def bst_to_dll(root):
    head = prev = [None, None]   # head[0], prev[0]
    def inorder(node):
        if not node: return
        inorder(node.left)
        if not prev[0]:
            head[0] = node
        else:
            prev[0].right = node
            node.left = prev[0]
        prev[0] = node
        inorder(node.right)
    inorder(root)
    return head[0]
\end{lstlisting}

\section{AVL Tree}

An AVL tree is a self-balancing BST where the heights of left and right subtrees differ by at most 1.
Four rotations maintain balance: LL, RR, LR, RL.

\begin{lstlisting}
class AVLNode:
    def __init__(self, val):
        self.val    = val
        self.left   = None
        self.right  = None
        self.height = 1

class AVLTree:
    def _h(self, node): return node.height if node else 0
    def _bf(self, node): return self._h(node.left) - self._h(node.right)
    def _update_h(self, node): node.height = 1 + max(self._h(node.left), self._h(node.right))

    def _rotate_right(self, y):
        x  = y.left; T2 = x.right
        x.right = y; y.left = T2
        self._update_h(y); self._update_h(x)
        return x

    def _rotate_left(self, x):
        y  = x.right; T2 = y.left
        y.left = x; x.right = T2
        self._update_h(x); self._update_h(y)
        return y

    def _rebalance(self, node):
        self._update_h(node)
        bf = self._bf(node)
        # LL
        if bf > 1 and self._bf(node.left) >= 0:
            return self._rotate_right(node)
        # LR
        if bf > 1 and self._bf(node.left) < 0:
            node.left = self._rotate_left(node.left)
            return self._rotate_right(node)
        # RR
        if bf < -1 and self._bf(node.right) <= 0:
            return self._rotate_left(node)
        # RL
        if bf < -1 and self._bf(node.right) > 0:
            node.right = self._rotate_right(node.right)
            return self._rotate_left(node)
        return node

    def insert(self, node, val):
        if not node: return AVLNode(val)
        if val < node.val:   node.left  = self.insert(node.left, val)
        elif val > node.val: node.right = self.insert(node.right, val)
        else: return node
        return self._rebalance(node)

    def delete(self, node, val):
        if not node: return None
        if val < node.val:   node.left  = self.delete(node.left, val)
        elif val > node.val: node.right = self.delete(node.right, val)
        else:
            if not node.left:  return node.right
            if not node.right: return node.left
            succ = node.right
            while succ.left: succ = succ.left
            node.val = succ.val
            node.right = self.delete(node.right, succ.val)
        return self._rebalance(node)
\end{lstlisting}

\begin{compbox}[BST / AVL Complexities]
\begin{tabular}{@{}llll@{}}
\toprule
\textbf{Operation} & \textbf{BST Avg} & \textbf{BST Worst} & \textbf{AVL} \\
\midrule
Search  & $O(\log n)$ & $O(n)$ & $O(\log n)$ \\
Insert  & $O(\log n)$ & $O(n)$ & $O(\log n)$ \\
Delete  & $O(\log n)$ & $O(n)$ & $O(\log n)$ \\
Space   & $O(n)$ & $O(n)$ & $O(n)$ \\
\bottomrule
\end{tabular}
\end{compbox}

\begin{tipbox}[Python Alternative to Balanced BST]
Python lacks a built-in balanced BST. Use \texttt{sortedcontainers.SortedList} from PyPI for $O(\log n)$ insert/delete/search on a sorted sequence. In competitive programming, use it in place of C++'s \texttt{std::set} or \texttt{std::multiset}.
\begin{lstlisting}
from sortedcontainers import SortedList
sl = SortedList()
sl.add(3); sl.add(1); sl.add(2)
sl.discard(1)
sl.bisect_left(2)    # index
\end{lstlisting}
\end{tipbox}

%%%==========================================================
\chapter{Heaps and Priority Queues}
%%%==========================================================

\section{Heap Operations via \texttt{heapq}}

\begin{lstlisting}
import heapq

# Python provides a MIN-HEAP only.
# For max-heap, negate all values.

# ── Basic min-heap ────────────────────────────────────────────
h = []
heapq.heappush(h, 5)
heapq.heappush(h, 2)
heapq.heappush(h, 8)
heapq.heappop(h)           # returns 2 (minimum)
h[0]                       # peek minimum without popping

# ── Heapify in O(n) ──────────────────────────────────────────
arr = [5, 3, 8, 1, 9]
heapq.heapify(arr)

# ── Max-heap via negation ─────────────────────────────────────
max_h = []
for x in [5, 3, 8, 1, 9]:
    heapq.heappush(max_h, -x)
-heapq.heappop(max_h)      # returns 9

# ── heappushpop / heapreplace ─────────────────────────────────
heapq.heappushpop(h, x)    # push x then pop min (more efficient)
heapq.heapreplace(h, x)    # pop min then push x (min <= x required)

# ── Custom objects: use (priority, counter, obj) ──────────────
import itertools
counter = itertools.count()   # tie-breaker for equal priorities

class PQ:
    def __init__(self):
        self._h = []
        self._counter = itertools.count()
        self._removed = set()

    def push(self, priority, item):
        cnt = next(self._counter)
        heapq.heappush(self._h, (priority, cnt, item))

    def pop(self):
        while self._h:
            priority, cnt, item = heapq.heappop(self._h)
            if cnt not in self._removed:
                return priority, item
        raise KeyError("pop from empty priority queue")

    def remove(self, cnt):
        self._removed.add(cnt)
\end{lstlisting}

\section{Heap Sort}

\begin{lstlisting}
# ── Heap Sort using heapq ── O(n log n)  SC: O(n) ─────────────
def heap_sort(arr):
    heapq.heapify(arr)
    return [heapq.heappop(arr) for _ in range(len(arr))]

# ── In-place Heap Sort ── O(n log n)  SC: O(1) ─────────────── 
def heap_sort_inplace(arr):
    n = len(arr)
    # Build max-heap
    for i in range(n // 2 - 1, -1, -1):
        _sift_down(arr, n, i)
    # Extract elements one by one
    for i in range(n-1, 0, -1):
        arr[0], arr[i] = arr[i], arr[0]
        _sift_down(arr, i, 0)

def _sift_down(arr, n, i):
    largest = i
    l, r = 2*i + 1, 2*i + 2
    if l < n and arr[l] > arr[largest]: largest = l
    if r < n and arr[r] > arr[largest]: largest = r
    if largest != i:
        arr[i], arr[largest] = arr[largest], arr[i]
        _sift_down(arr, n, largest)
\end{lstlisting}

\section{Heap Applications}

\begin{lstlisting}
# ── Top K Frequent Elements ── O(n log k) ─────────────────────
from collections import Counter
import heapq
def top_k_frequent(nums, k):
    freq = Counter(nums)
    return heapq.nlargest(k, freq.keys(), key=freq.get)

# ── K Closest Points to Origin ── O(n log k) ─────────────────
def k_closest(points, k):
    h = []
    for x, y in points:
        dist = x*x + y*y
        heapq.heappush(h, (-dist, x, y))
        if len(h) > k: heapq.heappop(h)
    return [[x, y] for _, x, y in h]

# ── Find Kth Smallest in Matrix (sorted rows & cols) ── O(k log n)
def kth_smallest_matrix(matrix, k):
    n = len(matrix)
    h = [(matrix[0][0], 0, 0)]
    visited = {(0, 0)}
    for _ in range(k - 1):
        val, r, c = heapq.heappop(h)
        for dr, dc in [(0, 1), (1, 0)]:
            nr, nc = r + dr, c + dc
            if nr < n and nc < n and (nr, nc) not in visited:
                heapq.heappush(h, (matrix[nr][nc], nr, nc))
                visited.add((nr, nc))
    return heapq.heappop(h)[0]

# ── Minimum Cost to Connect Ropes ── O(n log n) ───────────────
def min_cost_ropes(ropes):
    heapq.heapify(ropes)
    total = 0
    while len(ropes) > 1:
        a = heapq.heappop(ropes)
        b = heapq.heappop(ropes)
        cost = a + b
        total += cost
        heapq.heappush(ropes, cost)
    return total

# ── Sliding Window Median ── O(n log n) ───────────────────────
def sliding_window_median(nums, k):
    from sortedcontainers import SortedList
    window = SortedList()
    result = []
    for i, x in enumerate(nums):
        window.add(x)
        if len(window) > k:
            window.remove(nums[i - k])
        if len(window) == k:
            mid = k // 2
            if k % 2 == 1: result.append(float(window[mid]))
            else:          result.append((window[mid-1] + window[mid]) / 2.0)
    return result
\end{lstlisting}

%%%==========================================================
\chapter{Hash Tables and Sets}
%%%==========================================================

\section{Python dict}

\begin{lstlisting}
# ── dict (like unordered_map in C++) ─────────────────────────
d = {}
d = {'a': 1, 'b': 2}
d = dict(zip(keys, values))
d = dict.fromkeys(keys, 0)     # all values = 0

d['key'] = 'val'               # O(1) avg
d.get('key', default)          # safe get; no KeyError
d.pop('key')                   # remove and return
d.pop('key', None)             # safe pop
del d['key']
'key' in d                     # O(1) avg
d.keys(); d.values(); d.items()

# Merge dicts (Python 3.9+)
merged = d1 | d2               # new merged dict
d1 |= d2                       # update d1 with d2

# ── defaultdict ───────────────────────────────────────────────
from collections import defaultdict
graph   = defaultdict(list)      # missing key -> []
freq    = defaultdict(int)       # missing key -> 0
visited = defaultdict(bool)      # missing key -> False

# ── Counter ───────────────────────────────────────────────────
from collections import Counter
c = Counter([1,2,2,3,3,3])      # {3:3, 2:2, 1:1}
c.most_common(2)                 # [(3,3), (2,2)]
c['z']                           # 0 -- no KeyError
c.update([3, 4])                 # add more elements
c.subtract([1, 2])               # subtract counts
c1 + c2                          # sum (keeps positive)
c1 - c2                          # difference (drops non-positive)
c1 & c2                          # min of each
c1 | c2                          # max of each
+c                               # remove zero and negative counts

# ── OrderedDict ───────────────────────────────────────────────
from collections import OrderedDict
od = OrderedDict()
od.move_to_end('key')            # move to back (LRU cache)
od.move_to_end('key', last=False)# move to front
od.popitem(last=True)            # pop from back (LIFO)
od.popitem(last=False)           # pop from front (FIFO)
\end{lstlisting}

\section{Python set and frozenset}

\begin{lstlisting}
# ── set (like unordered_set in C++) ──────────────────────────
s = set()
s = {1, 2, 3}
s.add(4)                   # O(1)
s.remove(3)                # KeyError if absent
s.discard(3)               # safe remove
s.pop()                    # remove arbitrary element
3 in s                     # O(1)

# Set operations
s1 | s2                    # union
s1 & s2                    # intersection
s1 - s2                    # difference (in s1, not s2)
s1 ^ s2                    # symmetric difference
s1.issubset(s2)
s1.issuperset(s2)
s1.isdisjoint(s2)

# frozenset -- immutable, can be used as dict key
fs = frozenset([1, 2, 3])
\end{lstlisting}

\section{Hash Table Applications}

\begin{lstlisting}
# ── Two Sum ── O(n)  SC: O(n) ────────────────────────────────
def two_sum(nums, target):
    seen = {}
    for i, x in enumerate(nums):
        complement = target - x
        if complement in seen: return [seen[complement], i]
        seen[x] = i
    return []

# ── Group Anagrams ── O(n*k)  SC: O(n*k) ─────────────────────
def group_anagrams(strs):
    groups = defaultdict(list)
    for s in strs:
        key = tuple(sorted(s))
        groups[key].append(s)
    return list(groups.values())

# ── Longest Consecutive Sequence ── O(n) ─────────────────────
def longest_consecutive(nums):
    num_set = set(nums)
    best = 0
    for x in num_set:
        if x - 1 not in num_set:
            cur, length = x, 1
            while cur + 1 in num_set: cur += 1; length += 1
            best = max(best, length)
    return best

# ── Subarray with Zero Sum (exists?) ── O(n) ─────────────────
def has_zero_sum_subarray(arr):
    seen = {0}; prefix = 0
    for x in arr:
        prefix += x
        if prefix in seen: return True
        seen.add(prefix)
    return False

# ── Implement Custom HashMap ──────────────────────────────────
class MyHashMap:
    def __init__(self):
        self.SIZE = 1024
        self.buckets = [[] for _ in range(self.SIZE)]

    def _bucket(self, key): return key % self.SIZE

    def put(self, key, value):
        b = self._bucket(key)
        for i, (k, v) in enumerate(self.buckets[b]):
            if k == key: self.buckets[b][i] = (key, value); return
        self.buckets[b].append((key, value))

    def get(self, key):
        for k, v in self.buckets[self._bucket(key)]:
            if k == key: return v
        return -1

    def remove(self, key):
        b = self._bucket(key)
        self.buckets[b] = [(k, v) for k, v in self.buckets[b] if k != key]

# ── Isomorphic Strings ── O(n) ────────────────────────────────
def is_isomorphic(s, t):
    s2t = {}; t2s = {}
    for a, b in zip(s, t):
        if (a in s2t and s2t[a] != b) or (b in t2s and t2s[b] != a):
            return False
        s2t[a] = b; t2s[b] = a
    return True

# ── Word Pattern ── O(n) ──────────────────────────────────────
def word_pattern(pattern, s):
    words = s.split()
    if len(pattern) != len(words): return False
    p2w = {}; w2p = {}
    for p, w in zip(pattern, words):
        if (p in p2w and p2w[p] != w) or (w in w2p and w2p[w] != p):
            return False
        p2w[p] = w; w2p[w] = p
    return True
\end{lstlisting}

%%%==========================================================
\chapter{Tries}
%%%==========================================================

\section{Trie Implementation}

\begin{lstlisting}
class TrieNode:
    def __init__(self):
        self.children = {}    # char -> TrieNode
        self.is_end   = False
        self.count    = 0     # number of words through this node

class Trie:
    def __init__(self):
        self.root = TrieNode()

    # ── Insert ── O(m) where m = word length ─────────────────
    def insert(self, word):
        node = self.root
        for ch in word:
            if ch not in node.children:
                node.children[ch] = TrieNode()
            node = node.children[ch]
            node.count += 1
        node.is_end = True

    # ── Search exact word ── O(m) ────────────────────────────
    def search(self, word):
        node = self.root
        for ch in word:
            if ch not in node.children: return False
            node = node.children[ch]
        return node.is_end

    # ── Prefix exists ── O(m) ────────────────────────────────
    def starts_with(self, prefix):
        node = self.root
        for ch in prefix:
            if ch not in node.children: return False
            node = node.children[ch]
        return True

    # ── Delete ── O(m) ───────────────────────────────────────
    def delete(self, word):
        def _del(node, word, i):
            if i == len(word):
                if not node.is_end: return False
                node.is_end = False
                return len(node.children) == 0
            ch = word[i]
            if ch not in node.children: return False
            should_del = _del(node.children[ch], word, i+1)
            if should_del:
                node.children[ch].count -= 1
                if node.children[ch].count == 0:
                    del node.children[ch]
                    return not node.is_end and len(node.children) == 0
            return False
        _del(self.root, word, 0)

    # ── Count words with prefix ── O(m) ──────────────────────
    def count_with_prefix(self, prefix):
        node = self.root
        for ch in prefix:
            if ch not in node.children: return 0
            node = node.children[ch]
        return node.count

    # ── Longest Common Prefix of all words ── O(S) ───────────
    def longest_common_prefix(self):
        prefix = []
        node = self.root
        while len(node.children) == 1 and not node.is_end:
            ch = next(iter(node.children))
            prefix.append(ch)
            node = node.children[ch]
        return "".join(prefix)

    # ── Auto-complete: all words with prefix ── O(m + total)─
    def autocomplete(self, prefix):
        node = self.root
        for ch in prefix:
            if ch not in node.children: return []
            node = node.children[ch]
        result = []
        def dfs(n, curr):
            if n.is_end: result.append(prefix[:-len(curr)] + curr if curr else prefix)
            for ch, child in n.children.items():
                dfs(child, curr + ch)
        def dfs2(n, path):
            if n.is_end: result.append(path)
            for ch, child in n.children.items():
                dfs2(child, path + ch)
        dfs2(node, prefix)
        return result
\end{lstlisting}

\section{Trie Applications}

\begin{lstlisting}
# ── Word Search II (find all dictionary words in grid) ────────
def find_words(board, words):
    trie = Trie()
    for w in words: trie.insert(w)
    m, n = len(board), len(board[0])
    result = set()

    def dfs(node, r, c, path):
        ch = board[r][c]
        if ch not in node.children: return
        next_node = node.children[ch]
        path += ch
        if next_node.is_end: result.add(path)
        board[r][c] = '#'      # mark visited
        for dr, dc in [(-1,0),(1,0),(0,-1),(0,1)]:
            nr, nc = r+dr, c+dc
            if 0 <= nr < m and 0 <= nc < n and board[nr][nc] != '#':
                dfs(next_node, nr, nc, path)
        board[r][c] = ch       # restore

    for r in range(m):
        for c in range(n):
            dfs(trie.root, r, c, "")
    return list(result)

# ── Maximum XOR of Two Numbers using Trie ── O(n * 32) ────────
class BitTrie:
    def __init__(self): self.root = {}

    def insert(self, num):
        node = self.root
        for i in range(31, -1, -1):
            bit = (num >> i) & 1
            if bit not in node: node[bit] = {}
            node = node[bit]

    def max_xor(self, num):
        node = self.root
        xor = 0
        for i in range(31, -1, -1):
            bit = (num >> i) & 1
            want = 1 - bit
            if want in node:
                xor |= (1 << i); node = node[want]
            else:
                node = node.get(bit, {})
        return xor

def find_maximum_xor(nums):
    bt = BitTrie()
    for x in nums: bt.insert(x)
    return max(bt.max_xor(x) for x in nums)
\end{lstlisting}


%%%%%%%%%%%%%%%%%%%%%%%%%%%%%%%%%%%%%%%%%%%%%%%%%%%%%%%%%%%%%%%
\part{Advanced Data Structures}
%%%%%%%%%%%%%%%%%%%%%%%%%%%%%%%%%%%%%%%%%%%%%%%%%%%%%%%%%%%%%%%

%%%==========================================================
\chapter{Disjoint Set Union (DSU / Union-Find)}
%%%==========================================================

DSU efficiently answers \textbf{connectivity queries}: are two elements in the same set?
Two optimisations make every operation nearly $O(1)$: \textbf{path compression} and \textbf{union by rank/size}.

\section{Full DSU Implementation}

\begin{lstlisting}
class DSU:
    def __init__(self, n):
        self.parent = list(range(n))
        self.rank   = [0] * n
        self.size   = [1] * n
        self.num_components = n

    # Find with path compression -- O(alpha(n)) amortised
    def find(self, x):
        if self.parent[x] != x:
            self.parent[x] = self.find(self.parent[x])
        return self.parent[x]

    # Iterative path halving (avoids recursion-limit issues)
    def find_iter(self, x):
        while self.parent[x] != x:
            self.parent[x] = self.parent[self.parent[x]]   # path halving
            x = self.parent[x]
        return x

    # Union by rank -- O(alpha(n)) amortised
    def union(self, x, y):
        px, py = self.find(x), self.find(y)
        if px == py: return False    # already connected
        if self.rank[px] < self.rank[py]: px, py = py, px
        self.parent[py] = px
        self.size[px] += self.size[py]
        if self.rank[px] == self.rank[py]: self.rank[px] += 1
        self.num_components -= 1
        return True

    def connected(self, x, y): return self.find(x) == self.find(y)
    def get_size(self, x):      return self.size[self.find(x)]
    def components(self):       return self.num_components
\end{lstlisting}

\section{DSU Applications}

\begin{lstlisting}
from collections import defaultdict

# Detect Cycle in Undirected Graph -- O(E * alpha(V))
def has_cycle(n, edges):
    dsu = DSU(n)
    for u, v in edges:
        if not dsu.union(u, v): return True
    return False

# Kruskal's MST -- O(E log E)
def kruskal_mst(n, edges):
    edges.sort(key=lambda e: e[2])
    dsu = DSU(n)
    mst_cost = 0; mst_edges = []
    for u, v, w in edges:
        if dsu.union(u, v):
            mst_cost += w; mst_edges.append((u, v, w))
    return mst_cost, mst_edges

# Number of Connected Components
def count_components(n, edges):
    dsu = DSU(n)
    for u, v in edges: dsu.union(u, v)
    return dsu.components()

# Redundant Connection -- find edge that creates a cycle
def find_redundant_connection(edges):
    n = len(edges)
    dsu = DSU(n + 1)
    for u, v in edges:
        if not dsu.union(u, v): return [u, v]
    return []

# Accounts Merge -- group emails by owner
def accounts_merge(accounts):
    email_id = {}
    dsu = DSU(len(accounts))
    for i, acc in enumerate(accounts):
        for email in acc[1:]:
            if email in email_id: dsu.union(i, email_id[email])
            email_id[email] = i
    groups = defaultdict(set)
    for email, i in email_id.items():
        groups[dsu.find(i)].add(email)
    return [[accounts[i][0]] + sorted(emails) for i, emails in groups.items()]

# Minimum Number of Operations to Make Network Connected
def make_connected(n, connections):
    if len(connections) < n - 1: return -1   # not enough cables
    dsu = DSU(n)
    extra = 0
    for u, v in connections:
        if not dsu.union(u, v): extra += 1   # redundant edge
    components = dsu.components()
    return components - 1 if extra >= components - 1 else -1
\end{lstlisting}

%%%==========================================================
\chapter{Segment Trees}
%%%==========================================================

\section{Range Sum Segment Tree}

\begin{lstlisting}
class SegTree:
    def __init__(self, arr):
        self.n = len(arr)
        self.tree = [0] * (4 * self.n)
        if arr: self._build(arr, 0, 0, self.n - 1)

    def _build(self, arr, nd, lo, hi):
        if lo == hi: self.tree[nd] = arr[lo]; return
        mid = (lo + hi) // 2
        self._build(arr, 2*nd+1, lo, mid)
        self._build(arr, 2*nd+2, mid+1, hi)
        self.tree[nd] = self.tree[2*nd+1] + self.tree[2*nd+2]

    # Point Update -- O(log n)
    def update(self, idx, val, nd=0, lo=None, hi=None):
        if lo is None: lo, hi = 0, self.n - 1
        if lo == hi: self.tree[nd] = val; return
        mid = (lo + hi) // 2
        if idx <= mid: self.update(idx, val, 2*nd+1, lo, mid)
        else:          self.update(idx, val, 2*nd+2, mid+1, hi)
        self.tree[nd] = self.tree[2*nd+1] + self.tree[2*nd+2]

    # Range Sum Query -- O(log n)
    def query(self, l, r, nd=0, lo=None, hi=None):
        if lo is None: lo, hi = 0, self.n - 1
        if r < lo or hi < l: return 0
        if l <= lo and hi <= r: return self.tree[nd]
        mid = (lo + hi) // 2
        return (self.query(l, r, 2*nd+1, lo, mid) +
                self.query(l, r, 2*nd+2, mid+1, hi))
\end{lstlisting}

\section{Segment Tree with Lazy Propagation}

\begin{lstlisting}
class LazySegTree:
    # Range-add update on [l,r]; range sum query
    def __init__(self, arr):
        n = self.n = len(arr)
        self.tree = [0] * (4 * n)
        self.lazy = [0] * (4 * n)
        self._build(arr, 0, 0, n - 1)

    def _build(self, arr, nd, lo, hi):
        if lo == hi: self.tree[nd] = arr[lo]; return
        mid = (lo + hi) // 2
        self._build(arr, 2*nd+1, lo, mid)
        self._build(arr, 2*nd+2, mid+1, hi)
        self.tree[nd] = self.tree[2*nd+1] + self.tree[2*nd+2]

    def _push_down(self, nd, lo, hi):
        if not self.lazy[nd]: return
        mid = (lo + hi) // 2; d = self.lazy[nd]
        self.tree[2*nd+1] += d * (mid - lo + 1); self.lazy[2*nd+1] += d
        self.tree[2*nd+2] += d * (hi - mid);     self.lazy[2*nd+2] += d
        self.lazy[nd] = 0

    # Range add delta to arr[l..r] -- O(log n)
    def range_add(self, l, r, delta, nd=0, lo=None, hi=None):
        if lo is None: lo, hi = 0, self.n - 1
        if r < lo or hi < l: return
        if l <= lo and hi <= r:
            self.tree[nd] += delta * (hi - lo + 1); self.lazy[nd] += delta; return
        self._push_down(nd, lo, hi)
        mid = (lo + hi) // 2
        self.range_add(l, r, delta, 2*nd+1, lo, mid)
        self.range_add(l, r, delta, 2*nd+2, mid+1, hi)
        self.tree[nd] = self.tree[2*nd+1] + self.tree[2*nd+2]

    # Range sum query -- O(log n)
    def query(self, l, r, nd=0, lo=None, hi=None):
        if lo is None: lo, hi = 0, self.n - 1
        if r < lo or hi < l: return 0
        if l <= lo and hi <= r: return self.tree[nd]
        self._push_down(nd, lo, hi)
        mid = (lo + hi) // 2
        return (self.query(l,r,2*nd+1,lo,mid) + self.query(l,r,2*nd+2,mid+1,hi))
\end{lstlisting}

\begin{compbox}[Segment Tree Complexities]
\begin{tabular}{@{}lll@{}}
\toprule
\textbf{Operation} & \textbf{Time} & \textbf{Space} \\
\midrule
Build           & $O(n)$        & $O(4n)$ \\
Point update    & $O(\log n)$   & --- \\
Range query     & $O(\log n)$   & --- \\
Range update (lazy) & $O(\log n)$ & $O(4n)$ for lazy array \\
\bottomrule
\end{tabular}
\end{compbox}

%%%==========================================================
\chapter{Fenwick Tree (Binary Indexed Tree)}
%%%==========================================================

\begin{lstlisting}
class BIT:
    # 1-indexed BIT. Supports point update and prefix sum query.
    def __init__(self, n):
        self.n    = n
        self.tree = [0] * (n + 1)

    # Add delta to position i -- O(log n)
    def update(self, i, delta):
        while i <= self.n:
            self.tree[i] += delta
            i += i & (-i)      # add lowest set bit

    # Prefix sum [1..i] -- O(log n)
    def prefix(self, i):
        total = 0
        while i > 0:
            total += self.tree[i]
            i -= i & (-i)      # remove lowest set bit
        return total

    # Range sum [l..r]
    def range_sum(self, l, r): return self.prefix(r) - self.prefix(l - 1)

    # O(n) build from array (0-indexed input)
    def build(self, arr):
        for i, x in enumerate(arr, start=1):
            self.tree[i] += x
            j = i + (i & -i)
            if j <= self.n: self.tree[j] += self.tree[i]

    # Find kth element (1-indexed)
    def kth(self, k):
        idx = 0
        for pw in range(self.n.bit_length(), -1, -1):
            nxt = idx + (1 << pw)
            if nxt <= self.n and self.tree[nxt] < k:
                idx = nxt; k -= self.tree[nxt]
        return idx + 1


# Range Update, Point Query BIT
class BIT_RangeUpdate:
    # Add delta to range [l,r]; query single point i.
    def __init__(self, n): self.bit = BIT(n)

    def range_update(self, l, r, delta):
        self.bit.update(l, delta)
        self.bit.update(r + 1, -delta)

    def point_query(self, i): return self.bit.prefix(i)


# Range Update, Range Query using two BITs
class BIT_RangeRange:
    # Add delta to [l,r]; query sum of [l,r].
    def __init__(self, n):
        self.B1 = BIT(n); self.B2 = BIT(n); self.n = n

    def _prefix(self, i):
        return self.B1.prefix(i) * i - self.B2.prefix(i)

    def range_update(self, l, r, delta):
        self.B1.update(l, delta);       self.B1.update(r+1, -delta)
        self.B2.update(l, delta*(l-1)); self.B2.update(r+1, -delta*r)

    def range_query(self, l, r):
        return self._prefix(r) - self._prefix(l-1)


# 2D BIT -- O(log m * log n) per operation
class BIT2D:
    def __init__(self, m, n):
        self.m = m; self.n = n
        self.tree = [[0]*(n+1) for _ in range(m+1)]

    def update(self, r, c, delta):
        i = r
        while i <= self.m:
            j = c
            while j <= self.n:
                self.tree[i][j] += delta; j += j & (-j)
            i += i & (-i)

    def prefix(self, r, c):
        total = 0; i = r
        while i > 0:
            j = c
            while j > 0: total += self.tree[i][j]; j -= j & (-j)
            i -= i & (-i)
        return total

    def rect_query(self, r1, c1, r2, c2):
        return (self.prefix(r2,c2) - self.prefix(r1-1,c2)
                - self.prefix(r2,c1-1) + self.prefix(r1-1,c1-1))


# Count Inversions -- O(n log n)
def count_inversions(arr):
    if not arr: return 0
    compressed = {v: i+1 for i, v in enumerate(sorted(set(arr)))}
    bit = BIT(len(compressed))
    inversions = 0
    for x in arr:
        rank = compressed[x]
        inversions += bit.prefix(len(compressed)) - bit.prefix(rank)
        bit.update(rank, 1)
    return inversions
\end{lstlisting}

%%%==========================================================
\chapter{Sparse Table (Static RMQ)}
%%%==========================================================

A Sparse Table answers \textbf{Range Min/Max Queries} in $O(1)$ after $O(n \log n)$ preprocessing.
Requires an \textbf{idempotent} operation (min, max, gcd); does \emph{not} support updates.

\begin{lstlisting}
import math

class SparseTable:
    # Range Minimum Query in O(1) after O(n log n) build.
    def __init__(self, arr):
        n = len(arr)
        LOG = max(1, n.bit_length())
        self.log = [0] * (n + 1)
        for i in range(2, n + 1):
            self.log[i] = self.log[i // 2] + 1
        # st[k][i] = min of arr[i .. i + 2^k - 1]
        self.st = [arr[:]]
        for k in range(1, LOG):
            prev = self.st[k-1]
            half = 1 << (k-1)
            self.st.append([min(prev[i], prev[i + half])
                            for i in range(n - (1 << k) + 1)])

    # O(1) range minimum query
    def query(self, l, r):
        k = self.log[r - l + 1]
        return min(self.st[k][l], self.st[k][r - (1 << k) + 1])


class SparseTableMax:
    # Range Maximum Query in O(1).
    def __init__(self, arr):
        n = len(arr); LOG = max(1, n.bit_length())
        self.log = [0]*(n+1)
        for i in range(2,n+1): self.log[i] = self.log[i//2]+1
        self.st = [arr[:]]
        for k in range(1, LOG):
            prev = self.st[k-1]; half = 1 << (k-1)
            self.st.append([max(prev[i], prev[i+half])
                            for i in range(n-(1<<k)+1)])

    def query(self, l, r):
        k = self.log[r-l+1]
        return max(self.st[k][l], self.st[k][r-(1<<k)+1])


# GCD Sparse Table
class GCDSparseTable:
    def __init__(self, arr):
        import math
        n = len(arr); LOG = max(1, n.bit_length())
        self.log = [0]*(n+1)
        for i in range(2,n+1): self.log[i]=self.log[i//2]+1
        self.st = [arr[:]]
        for k in range(1, LOG):
            prev = self.st[k-1]; half = 1 << (k-1)
            self.st.append([math.gcd(prev[i], prev[i+half])
                            for i in range(n-(1<<k)+1)])

    def query(self, l, r):
        k = self.log[r-l+1]
        return math.gcd(self.st[k][l], self.st[k][r-(1<<k)+1])
\end{lstlisting}

\begin{compbox}[Sparse Table vs Segment Tree vs BIT]
\begin{tabular}{@{}llll@{}}
\toprule
\textbf{Structure} & \textbf{Build} & \textbf{Query} & \textbf{Update} \\
\midrule
Sparse Table     & $O(n \log n)$ & $O(1)$        & None \\
Segment Tree     & $O(n)$        & $O(\log n)$   & $O(\log n)$ \\
BIT              & $O(n)$        & $O(\log n)$   & $O(\log n)$ \\
\bottomrule
\end{tabular}\\[4pt]
Use Sparse Table for static arrays needing ultra-fast idempotent range queries.
\end{compbox}


%%%%%%%%%%%%%%%%%%%%%%%%%%%%%%%%%%%%%%%%%%%%%%%%%%%%%%%%%%%%%%%
\part{Graph Algorithms}
%%%%%%%%%%%%%%%%%%%%%%%%%%%%%%%%%%%%%%%%%%%%%%%%%%%%%%%%%%%%%%%

%%%==========================================================
\chapter{Graph Representations and Traversals}
%%%==========================================================

\section{Graph Representations}

\begin{lstlisting}
from collections import defaultdict, deque

# -- Adjacency List (most common) -- Space: O(V + E)
graph = defaultdict(list)
# Undirected
graph[0].append(1); graph[1].append(0)
# Directed weighted
graph[u].append((v, w))

# -- Adjacency Matrix -- Space: O(V^2)
n = 5
adj = [[0]*n for _ in range(n)]
adj[u][v] = 1           # unweighted
adj[u][v] = w           # weighted

# -- Edge List
edges = [(0,1,5), (1,2,3), ...]    # (u, v, weight)

# -- Building from input
def build_graph(n, edge_list, directed=False):
    g = defaultdict(list)
    for u, v, *w in edge_list:
        wt = w[0] if w else 1
        g[u].append((v, wt))
        if not directed: g[v].append((u, wt))
    return g
\end{lstlisting}

\section{Depth-First Search (DFS)}

\begin{lstlisting}
# -- Recursive DFS -- O(V + E)
def dfs_recursive(graph, start):
    visited = set()
    order   = []
    def dfs(u):
        visited.add(u)
        order.append(u)
        for v in graph[u]:
            if v not in visited:
                dfs(v)
    dfs(start)
    return order

# -- Iterative DFS (avoids recursion limit on large graphs)
def dfs_iterative(graph, start):
    visited = set([start])
    stack   = [start]
    order   = []
    while stack:
        u = stack.pop()
        order.append(u)
        for v in graph[u]:
            if v not in visited:
                visited.add(v)
                stack.append(v)
    return order

# -- DFS on entire graph (handles disconnected components)
def dfs_all(graph, n):
    visited = set()
    components = []
    for src in range(n):
        if src not in visited:
            comp = []
            stack = [src]
            visited.add(src)
            while stack:
                u = stack.pop()
                comp.append(u)
                for v in graph[u]:
                    if v not in visited:
                        visited.add(v)
                        stack.append(v)
            components.append(comp)
    return components
\end{lstlisting}

\section{Breadth-First Search (BFS)}

\begin{lstlisting}
# -- BFS -- O(V + E), finds shortest path in unweighted graph
def bfs(graph, start):
    visited = set([start])
    queue   = deque([start])
    order   = []
    while queue:
        u = queue.popleft()
        order.append(u)
        for v in graph[u]:
            if v not in visited:
                visited.add(v)
                queue.append(v)
    return order

# -- BFS with shortest distance tracking
def bfs_distance(graph, start):
    dist    = {start: 0}
    queue   = deque([start])
    while queue:
        u = queue.popleft()
        for v in graph[u]:
            if v not in dist:
                dist[v] = dist[u] + 1
                queue.append(v)
    return dist

# -- BFS with path reconstruction
def bfs_path(graph, start, end):
    parent  = {start: None}
    queue   = deque([start])
    while queue:
        u = queue.popleft()
        if u == end:
            path = []
            while u is not None:
                path.append(u); u = parent[u]
            return path[::-1]
        for v in graph[u]:
            if v not in parent:
                parent[v] = u
                queue.append(v)
    return []    # no path

# -- Multi-source BFS
def bfs_multi_source(grid, sources):
    m, n = len(grid), len(grid[0])
    dist  = [[-1]*n for _ in range(m)]
    queue = deque()
    for r, c in sources:
        dist[r][c] = 0
        queue.append((r, c))
    while queue:
        r, c = queue.popleft()
        for dr, dc in [(-1,0),(1,0),(0,-1),(0,1)]:
            nr, nc = r+dr, c+dc
            if 0 <= nr < m and 0 <= nc < n and dist[nr][nc] == -1:
                dist[nr][nc] = dist[r][c] + 1
                queue.append((nr, nc))
    return dist
\end{lstlisting}

\section{Topological Sort}

\begin{lstlisting}
# -- Kahn's Algorithm (BFS-based) -- O(V + E)
def topo_sort_kahn(graph, n):
    in_degree = [0] * n
    for u in range(n):
        for v in graph[u]: in_degree[v] += 1
    queue = deque(u for u in range(n) if in_degree[u] == 0)
    order = []
    while queue:
        u = queue.popleft()
        order.append(u)
        for v in graph[u]:
            in_degree[v] -= 1
            if in_degree[v] == 0: queue.append(v)
    return order if len(order) == n else []  # empty => cycle

# -- DFS-based Topological Sort -- O(V + E)
def topo_sort_dfs(graph, n):
    WHITE, GRAY, BLACK = 0, 1, 2
    color = [WHITE] * n
    order = []
    has_cycle = [False]

    def dfs(u):
        color[u] = GRAY
        for v in graph[u]:
            if color[v] == GRAY:
                has_cycle[0] = True; return
            if color[v] == WHITE: dfs(v)
        color[u] = BLACK
        order.append(u)

    for u in range(n):
        if color[u] == WHITE: dfs(u)
    return [] if has_cycle[0] else order[::-1]

# -- Course Schedule (detect cycle via topo sort) -- O(V+E)
def can_finish(num_courses, prerequisites):
    graph = defaultdict(list)
    for a, b in prerequisites: graph[b].append(a)
    return len(topo_sort_kahn(graph, num_courses)) == num_courses
\end{lstlisting}

%%%==========================================================
\chapter{Shortest Path Algorithms}
%%%==========================================================

\section{Dijkstra's Algorithm}

\begin{lstlisting}
import heapq

# -- Dijkstra -- O((V + E) log V) with binary heap
def dijkstra(graph, src, n):
    # graph[u] = [(v, w), ...]
    dist = [float('inf')] * n
    dist[src] = 0
    heap = [(0, src)]    # (dist, node)
    while heap:
        d, u = heapq.heappop(heap)
        if d > dist[u]: continue    # stale entry
        for v, w in graph[u]:
            if dist[u] + w < dist[v]:
                dist[v] = dist[u] + w
                heapq.heappush(heap, (dist[v], v))
    return dist

# -- Dijkstra with path reconstruction
def dijkstra_path(graph, src, dst, n):
    dist   = [float('inf')] * n
    parent = [-1] * n
    dist[src] = 0
    heap = [(0, src)]
    while heap:
        d, u = heapq.heappop(heap)
        if d > dist[u]: continue
        if u == dst: break
        for v, w in graph[u]:
            if dist[u] + w < dist[v]:
                dist[v] = dist[u] + w
                parent[v] = u
                heapq.heappush(heap, (dist[v], v))
    if dist[dst] == float('inf'): return [], float('inf')
    path = []
    cur  = dst
    while cur != -1: path.append(cur); cur = parent[cur]
    return path[::-1], dist[dst]

# -- Dijkstra on grid
def dijkstra_grid(grid, start, end):
    m, n = len(grid), len(grid[0])
    dist = [[float('inf')]*n for _ in range(m)]
    sr, sc = start; er, ec = end
    dist[sr][sc] = 0
    heap = [(0, sr, sc)]
    while heap:
        d, r, c = heapq.heappop(heap)
        if d > dist[r][c]: continue
        if (r, c) == (er, ec): return d
        for dr, dc in [(-1,0),(1,0),(0,-1),(0,1)]:
            nr, nc = r+dr, c+dc
            if 0 <= nr < m and 0 <= nc < n:
                nd = d + grid[nr][nc]
                if nd < dist[nr][nc]:
                    dist[nr][nc] = nd
                    heapq.heappush(heap, (nd, nr, nc))
    return dist[er][ec]
\end{lstlisting}

\section{Bellman-Ford Algorithm}

\begin{lstlisting}
# -- Bellman-Ford -- O(V * E), handles negative weights
def bellman_ford(n, edges, src):
    # edges: [(u, v, w), ...]
    dist = [float('inf')] * n
    dist[src] = 0
    for _ in range(n - 1):       # relax all edges n-1 times
        updated = False
        for u, v, w in edges:
            if dist[u] != float('inf') and dist[u] + w < dist[v]:
                dist[v] = dist[u] + w
                updated = True
        if not updated: break    # early termination

    # Detect negative cycle
    for u, v, w in edges:
        if dist[u] != float('inf') and dist[u] + w < dist[v]:
            return None          # negative cycle exists
    return dist

# -- SPFA (queue-based Bellman-Ford, faster in practice)
def spfa(n, graph, src):
    dist    = [float('inf')] * n
    in_queue = [False] * n
    cnt     = [0] * n           # count relaxations
    dist[src] = 0
    queue = deque([src])
    in_queue[src] = True
    while queue:
        u = queue.popleft()
        in_queue[u] = False
        for v, w in graph[u]:
            if dist[u] + w < dist[v]:
                dist[v] = dist[u] + w
                cnt[v] += 1
                if cnt[v] >= n: return None   # negative cycle
                if not in_queue[v]:
                    queue.append(v); in_queue[v] = True
    return dist
\end{lstlisting}

\section{Floyd-Warshall Algorithm}

\begin{lstlisting}
# -- Floyd-Warshall -- O(V^3), all-pairs shortest paths
def floyd_warshall(n, edges):
    INF = float('inf')
    dist = [[INF]*n for _ in range(n)]
    for i in range(n): dist[i][i] = 0
    for u, v, w in edges:
        dist[u][v] = min(dist[u][v], w)
        dist[v][u] = min(dist[v][u], w)    # remove for directed

    for k in range(n):
        for i in range(n):
            for j in range(n):
                if dist[i][k] + dist[k][j] < dist[i][j]:
                    dist[i][j] = dist[i][k] + dist[k][j]

    # Check for negative cycles (dist[i][i] < 0)
    for i in range(n):
        if dist[i][i] < 0: return None    # negative cycle

    return dist

# -- Reconstruct path from Floyd-Warshall
def floyd_warshall_path(n, edges):
    INF = float('inf')
    dist = [[INF]*n for _ in range(n)]
    next_node = [[None]*n for _ in range(n)]
    for i in range(n): dist[i][i] = 0
    for u, v, w in edges:
        if w < dist[u][v]: dist[u][v] = w; next_node[u][v] = v
    for k in range(n):
        for i in range(n):
            for j in range(n):
                if dist[i][k] + dist[k][j] < dist[i][j]:
                    dist[i][j] = dist[i][k] + dist[k][j]
                    next_node[i][j] = next_node[i][k]
    def get_path(u, v):
        if next_node[u][v] is None: return []
        path = [u]
        while u != v: u = next_node[u][v]; path.append(u)
        return path
    return dist, get_path
\end{lstlisting}

\begin{compbox}[Shortest Path Algorithm Comparison]
\begin{tabular}{@{}lllll@{}}
\toprule
\textbf{Algorithm} & \textbf{Time} & \textbf{Negative Weights} & \textbf{Neg. Cycles} & \textbf{Use When} \\
\midrule
BFS              & $O(V+E)$      & No           & No   & Unweighted graph \\
Dijkstra         & $O((V+E)\log V)$ & No        & No   & Non-negative weights \\
Bellman-Ford     & $O(VE)$       & Yes          & Detect & Negative weights \\
Floyd-Warshall   & $O(V^3)$      & Yes          & Detect & All-pairs, small $V$ \\
\bottomrule
\end{tabular}
\end{compbox}

%%%==========================================================
\chapter{Minimum Spanning Trees}
%%%==========================================================

\section{Kruskal's Algorithm}

\begin{lstlisting}
# -- Kruskal's MST -- O(E log E)
def kruskal(n, edges):
    # edges: [(w, u, v), ...]
    edges.sort()
    dsu = DSU(n)
    mst_weight = 0
    mst_edges  = []
    for w, u, v in edges:
        if dsu.union(u, v):
            mst_weight += w
            mst_edges.append((u, v, w))
            if len(mst_edges) == n - 1: break
    return mst_weight, mst_edges
\end{lstlisting}

\section{Prim's Algorithm}

\begin{lstlisting}
# -- Prim's MST (heap-based) -- O((V + E) log V)
def prim(graph, n):
    # graph[u] = [(v, w), ...]
    in_mst  = [False] * n
    key     = [float('inf')] * n
    parent  = [-1] * n
    key[0]  = 0
    heap    = [(0, 0)]    # (weight, vertex)
    mst_cost = 0

    while heap:
        w, u = heapq.heappop(heap)
        if in_mst[u]: continue
        in_mst[u]  = True
        mst_cost  += w
        for v, wt in graph[u]:
            if not in_mst[v] and wt < key[v]:
                key[v]    = wt
                parent[v] = u
                heapq.heappush(heap, (wt, v))

    return mst_cost, [(parent[i], i, key[i]) for i in range(1, n)]
\end{lstlisting}

%%%==========================================================
\chapter{Cycle Detection and Graph Properties}
%%%==========================================================

\section{Cycle Detection}

\begin{lstlisting}
# -- Cycle in Directed Graph (DFS coloring) -- O(V + E)
def has_cycle_directed(graph, n):
    WHITE, GRAY, BLACK = 0, 1, 2
    color = [WHITE] * n

    def dfs(u):
        color[u] = GRAY
        for v in graph[u]:
            if color[v] == GRAY: return True    # back edge
            if color[v] == WHITE and dfs(v): return True
        color[u] = BLACK
        return False

    return any(color[u] == WHITE and dfs(u) for u in range(n))

# -- Cycle in Undirected Graph (DFS) -- O(V + E)
def has_cycle_undirected(graph, n):
    visited = [False] * n

    def dfs(u, parent):
        visited[u] = True
        for v in graph[u]:
            if not visited[v]:
                if dfs(v, u): return True
            elif v != parent:
                return True              # back edge to non-parent
        return False

    return any(not visited[u] and dfs(u, -1) for u in range(n))
\end{lstlisting}

\section{Strongly Connected Components (SCCs)}

\begin{lstlisting}
# -- Kosaraju's Algorithm -- O(V + E)
def kosaraju(graph, n):
    # graph: adjacency list (0-indexed)
    visited = [False] * n
    order   = []

    def dfs1(u):
        visited[u] = True
        for v in graph[u]:
            if not visited[v]: dfs1(v)
        order.append(u)

    for u in range(n):
        if not visited[u]: dfs1(u)

    # Build reverse graph
    rev = defaultdict(list)
    for u in range(n):
        for v in graph[u]: rev[v].append(u)

    visited = [False] * n
    sccs    = []

    def dfs2(u, comp):
        visited[u] = True; comp.append(u)
        for v in rev[u]:
            if not visited[v]: dfs2(v, comp)

    for u in reversed(order):
        if not visited[u]:
            comp = []; dfs2(u, comp); sccs.append(comp)

    return sccs

# -- Tarjan's Algorithm -- O(V + E), single DFS
def tarjan_scc(graph, n):
    index    = [0]
    idx      = [-1] * n
    low      = [0]  * n
    on_stack = [False] * n
    stack    = []
    sccs     = []

    def dfs(u):
        idx[u] = low[u] = index[0]; index[0] += 1
        stack.append(u); on_stack[u] = True

        for v in graph[u]:
            if idx[v] == -1:
                dfs(v); low[u] = min(low[u], low[v])
            elif on_stack[v]:
                low[u] = min(low[u], idx[v])

        if low[u] == idx[u]:   # root of SCC
            scc = []
            while True:
                w = stack.pop(); on_stack[w] = False
                scc.append(w)
                if w == u: break
            sccs.append(scc)

    for u in range(n):
        if idx[u] == -1: dfs(u)

    return sccs
\end{lstlisting}

\section{Bridges and Articulation Points}

\begin{lstlisting}
# -- Find Bridges (edges whose removal disconnects graph) -- O(V+E)
def find_bridges(graph, n):
    disc   = [-1] * n
    low    = [0]  * n
    timer  = [0]
    bridges = []

    def dfs(u, parent):
        disc[u] = low[u] = timer[0]; timer[0] += 1
        for v in graph[u]:
            if v == parent: continue
            if disc[v] == -1:
                dfs(v, u)
                low[u] = min(low[u], low[v])
                if low[v] > disc[u]:
                    bridges.append((u, v))
            else:
                low[u] = min(low[u], disc[v])

    for u in range(n):
        if disc[u] == -1: dfs(u, -1)
    return bridges

# -- Find Articulation Points -- O(V + E)
def find_articulation_points(graph, n):
    disc     = [-1] * n
    low      = [0]  * n
    timer    = [0]
    is_ap    = [False] * n

    def dfs(u, parent):
        disc[u] = low[u] = timer[0]; timer[0] += 1
        children = 0
        for v in graph[u]:
            if v == parent: continue
            if disc[v] == -1:
                children += 1
                dfs(v, u)
                low[u] = min(low[u], low[v])
                # u is AP if: root with 2+ children, or non-root with low[v] >= disc[u]
                if parent == -1 and children > 1: is_ap[u] = True
                if parent != -1 and low[v] >= disc[u]: is_ap[u] = True
            else:
                low[u] = min(low[u], disc[v])

    for u in range(n):
        if disc[u] == -1: dfs(u, -1)
    return [u for u in range(n) if is_ap[u]]
\end{lstlisting}

\section{Bipartite Check}

\begin{lstlisting}
# -- Bipartite Check (2-coloring via BFS) -- O(V + E)
def is_bipartite(graph, n):
    color = [-1] * n
    for start in range(n):
        if color[start] != -1: continue
        queue = deque([start])
        color[start] = 0
        while queue:
            u = queue.popleft()
            for v in graph[u]:
                if color[v] == -1:
                    color[v] = 1 - color[u]
                    queue.append(v)
                elif color[v] == color[u]:
                    return False
    return True
\end{lstlisting}

%%%==========================================================
\chapter{Advanced Graph Algorithms}
%%%==========================================================

\section{Flood Fill / Island Problems}

\begin{lstlisting}
# -- Number of Islands (DFS) -- O(m*n)
def num_islands(grid):
    if not grid: return 0
    m, n   = len(grid), len(grid[0])
    count  = 0

    def dfs(r, c):
        if r < 0 or r >= m or c < 0 or c >= n or grid[r][c] != '1': return
        grid[r][c] = '0'    # mark visited (in-place)
        for dr, dc in [(-1,0),(1,0),(0,-1),(0,1)]: dfs(r+dr, c+dc)

    for r in range(m):
        for c in range(n):
            if grid[r][c] == '1': count += 1; dfs(r, c)
    return count

# -- Surrounded Regions -- O(m*n)
def solve_surrounded(board):
    if not board: return
    m, n = len(board), len(board[0])

    def dfs(r, c):
        if r < 0 or r >= m or c < 0 or c >= n or board[r][c] != 'O': return
        board[r][c] = 'S'    # safe
        for dr, dc in [(-1,0),(1,0),(0,-1),(0,1)]: dfs(r+dr, c+dc)

    # Mark border-connected O's as safe
    for r in range(m):
        if board[r][0] == 'O': dfs(r, 0)
        if board[r][n-1] == 'O': dfs(r, n-1)
    for c in range(n):
        if board[0][c] == 'O': dfs(0, c)
        if board[m-1][c] == 'O': dfs(m-1, c)

    for r in range(m):
        for c in range(n):
            if   board[r][c] == 'O': board[r][c] = 'X'
            elif board[r][c] == 'S': board[r][c] = 'O'
\end{lstlisting}

\section{0-1 BFS}

\begin{lstlisting}
# -- 0-1 BFS: edges have weight 0 or 1 -- O(V + E)
# Use deque: append front for weight-0, append back for weight-1
def bfs_01(graph, src, n):
    dist  = [float('inf')] * n
    dist[src] = 0
    dq    = deque([src])
    while dq:
        u = dq.popleft()
        for v, w in graph[u]:
            if dist[u] + w < dist[v]:
                dist[v] = dist[u] + w
                if w == 0: dq.appendleft(v)
                else:      dq.append(v)
    return dist
\end{lstlisting}

\section{Euler Path and Circuit}

\begin{lstlisting}
# -- Hierholzer's Algorithm -- Eulerian Circuit -- O(E)
def eulerian_circuit(graph, n):
    # Precondition: all vertices have even degree
    # graph is adjacency list with edge counts (multigraph)
    adj   = [list(graph[u]) for u in range(n)]
    stack = [0]
    path  = []
    ptr   = [0] * n    # iterator pointer (avoid O(n) per removal)
    # Use iterator approach for efficiency
    adj_iter = [iter(graph[u]) for u in range(n)]
    # Simpler version with sets:
    edge_used = set()
    adj2 = {u: list(v for v in graph[u]) for u in range(n)}
    stack2 = [0]; circuit = []
    while stack2:
        v = stack2[-1]
        if adj2[v]:
            u = adj2[v].pop()
            adj2[u].remove(v)   # undirected: remove back edge
            stack2.append(u)
        else:
            circuit.append(stack2.pop())
    return circuit[::-1]

# -- Check Eulerian conditions
def has_eulerian_circuit(graph, n):
    return all(len(graph[u]) % 2 == 0 for u in range(n))

def has_eulerian_path(graph, n):
    odd = sum(1 for u in range(n) if len(graph[u]) % 2 == 1)
    return odd == 0 or odd == 2
\end{lstlisting}

\section{Bellman-Ford for Cheapest Flights within K Stops}

\begin{lstlisting}
# -- Cheapest flights within k stops -- O(k * E)
def find_cheapest_price(n, flights, src, dst, k):
    INF  = float('inf')
    dist = [INF] * n
    dist[src] = 0
    for _ in range(k + 1):
        tmp = dist[:]
        for u, v, w in flights:
            if dist[u] != INF and dist[u] + w < tmp[v]:
                tmp[v] = dist[u] + w
        dist = tmp
    return dist[dst] if dist[dst] != INF else -1
\end{lstlisting}


%%%%%%%%%%%%%%%%%%%%%%%%%%%%%%%%%%%%%%%%%%%%%%%%%%%%%%%%%%%%%%%
\part{Sorting and Searching}
%%%%%%%%%%%%%%%%%%%%%%%%%%%%%%%%%%%%%%%%%%%%%%%%%%%%%%%%%%%%%%%

%%%==========================================================
\chapter{Sorting Algorithms}
%%%==========================================================

\begin{notebox}[Python's Built-in Sort]
\texttt{list.sort()} and \texttt{sorted()} use \textbf{Timsort} --- $O(n \log n)$ worst case, $O(n)$ on nearly-sorted data. Always prefer these in practice. Implement manual sorts only when the algorithm itself is being tested.
\end{notebox}

\section{Comparison-Based Sorts}

\begin{lstlisting}
# ============================================================
# BUBBLE SORT -- O(n^2) avg/worst, O(n) best, Stable
# ============================================================
def bubble_sort(arr):
    n = len(arr)
    for i in range(n):
        swapped = False
        for j in range(n - i - 1):
            if arr[j] > arr[j+1]:
                arr[j], arr[j+1] = arr[j+1], arr[j]
                swapped = True
        if not swapped: break    # already sorted

# ============================================================
# SELECTION SORT -- O(n^2), Not stable
# ============================================================
def selection_sort(arr):
    n = len(arr)
    for i in range(n):
        min_idx = i
        for j in range(i+1, n):
            if arr[j] < arr[min_idx]: min_idx = j
        arr[i], arr[min_idx] = arr[min_idx], arr[i]

# ============================================================
# INSERTION SORT -- O(n^2) avg/worst, O(n) best, Stable
# ============================================================
def insertion_sort(arr):
    for i in range(1, len(arr)):
        key = arr[i]; j = i - 1
        while j >= 0 and arr[j] > key:
            arr[j+1] = arr[j]; j -= 1
        arr[j+1] = key

# ============================================================
# MERGE SORT -- O(n log n) all cases, Stable, O(n) space
# ============================================================
def merge_sort(arr):
    if len(arr) <= 1: return arr
    mid   = len(arr) // 2
    left  = merge_sort(arr[:mid])
    right = merge_sort(arr[mid:])
    return _merge(left, right)

def _merge(left, right):
    result = []; i = j = 0
    while i < len(left) and j < len(right):
        if left[i] <= right[j]: result.append(left[i]); i += 1
        else:                    result.append(right[j]); j += 1
    result.extend(left[i:]); result.extend(right[j:])
    return result

# Bottom-up (iterative) merge sort
def merge_sort_iterative(arr):
    n = len(arr)
    width = 1
    while width < n:
        for lo in range(0, n, 2*width):
            mid  = min(lo + width, n)
            hi   = min(lo + 2*width, n)
            arr[lo:hi] = _merge(arr[lo:mid], arr[mid:hi])
        width *= 2
    return arr

# ============================================================
# QUICK SORT -- O(n log n) avg, O(n^2) worst, Not stable
# ============================================================
def quick_sort(arr, lo=0, hi=None):
    if hi is None: hi = len(arr) - 1
    if lo >= hi: return
    p = _partition(arr, lo, hi)
    quick_sort(arr, lo, p - 1)
    quick_sort(arr, p + 1, hi)

def _partition(arr, lo, hi):
    # Median-of-three pivot
    mid = (lo + hi) // 2
    if arr[mid] < arr[lo]: arr[lo], arr[mid] = arr[mid], arr[lo]
    if arr[hi]  < arr[lo]: arr[lo], arr[hi]  = arr[hi],  arr[lo]
    if arr[mid] < arr[hi]: arr[mid], arr[hi] = arr[hi],  arr[mid]
    pivot = arr[hi]
    i = lo - 1
    for j in range(lo, hi):
        if arr[j] <= pivot:
            i += 1; arr[i], arr[j] = arr[j], arr[i]
    arr[i+1], arr[hi] = arr[hi], arr[i+1]
    return i + 1

# Randomised quick sort
import random
def quick_sort_rand(arr, lo=0, hi=None):
    if hi is None: hi = len(arr) - 1
    if lo >= hi: return
    rand_idx = random.randint(lo, hi)
    arr[rand_idx], arr[hi] = arr[hi], arr[rand_idx]
    p = _partition(arr, lo, hi)
    quick_sort_rand(arr, lo, p - 1)
    quick_sort_rand(arr, p + 1, hi)

# 3-way partition (Dutch National Flag) for many duplicates
def quick_sort_3way(arr, lo=0, hi=None):
    if hi is None: hi = len(arr) - 1
    if lo >= hi: return
    lt, gt = lo, hi
    pivot  = arr[lo]; i = lo + 1
    while i <= gt:
        if   arr[i] < pivot: arr[lt], arr[i] = arr[i], arr[lt]; lt += 1; i += 1
        elif arr[i] > pivot: arr[i], arr[gt]  = arr[gt], arr[i]; gt -= 1
        else: i += 1
    quick_sort_3way(arr, lo, lt - 1)
    quick_sort_3way(arr, gt + 1, hi)

# ============================================================
# HEAP SORT -- O(n log n) all cases, Not stable, O(1) space
# ============================================================
def heap_sort(arr):
    n = len(arr)
    for i in range(n//2 - 1, -1, -1):
        _sift_down(arr, n, i)
    for i in range(n-1, 0, -1):
        arr[0], arr[i] = arr[i], arr[0]
        _sift_down(arr, i, 0)

def _sift_down(arr, n, i):
    largest = i; l = 2*i+1; r = 2*i+2
    if l < n and arr[l] > arr[largest]: largest = l
    if r < n and arr[r] > arr[largest]: largest = r
    if largest != i:
        arr[i], arr[largest] = arr[largest], arr[i]
        _sift_down(arr, n, largest)

# ============================================================
# SHELL SORT -- O(n log^2 n) with Ciura gaps, Not stable
# ============================================================
def shell_sort(arr):
    gaps = [701, 301, 132, 57, 23, 10, 4, 1]   # Ciura gap sequence
    for gap in gaps:
        for i in range(gap, len(arr)):
            temp = arr[i]; j = i
            while j >= gap and arr[j-gap] > temp:
                arr[j] = arr[j-gap]; j -= gap
            arr[j] = temp
\end{lstlisting}

\section{Non-Comparison Sorts}

\begin{lstlisting}
# ============================================================
# COUNTING SORT -- O(n + k), Stable, for integers in [0, k]
# ============================================================
def counting_sort(arr, k=None):
    if not arr: return arr
    if k is None: k = max(arr)
    count = [0] * (k + 1)
    for x in arr: count[x] += 1
    for i in range(1, k+1): count[i] += count[i-1]   # prefix sums
    output = [0] * len(arr)
    for x in reversed(arr):    # reversed for stability
        count[x] -= 1
        output[count[x]] = x
    return output

# ============================================================
# RADIX SORT -- O(d * (n + b)), Stable, d=digits, b=base
# ============================================================
def radix_sort(arr, base=10):
    if not arr: return arr
    max_val = max(arr)
    exp = 1
    while max_val // exp > 0:
        arr = _counting_sort_by_digit(arr, exp, base)
        exp *= base
    return arr

def _counting_sort_by_digit(arr, exp, base):
    n = len(arr)
    count  = [0] * base
    output = [0] * n
    for x in arr: count[(x // exp) % base] += 1
    for i in range(1, base): count[i] += count[i-1]
    for x in reversed(arr):
        idx = (x // exp) % base
        count[idx] -= 1
        output[count[idx]] = x
    return output

# ============================================================
# BUCKET SORT -- O(n + k) avg for uniformly distributed floats
# ============================================================
def bucket_sort(arr):
    if not arr: return arr
    n       = len(arr)
    buckets = [[] for _ in range(n)]
    min_v, max_v = min(arr), max(arr)
    span = max_v - min_v + 1e-9
    for x in arr:
        idx = int((x - min_v) / span * (n - 1))
        buckets[idx].append(x)
    result = []
    for b in buckets:
        b.sort()
        result.extend(b)
    return result
\end{lstlisting}

\section{Python Sort Customisation}

\begin{lstlisting}
# -- Custom key functions
arr.sort(key=lambda x: x[1])              # sort by 2nd element
arr.sort(key=lambda x: (-x[0], x[1]))     # desc by first, asc by second
arr.sort(key=len)                          # sort strings by length
arr.sort(key=str.lower)                    # case-insensitive

# -- functools.cmp_to_key (for custom comparator like C++ cmp)
from functools import cmp_to_key
def my_cmp(a, b):
    return -1 if str(a)+str(b) > str(b)+str(a) else 1  # largest number
nums = [3, 30, 34, 5, 9]
nums.sort(key=cmp_to_key(my_cmp))
print("".join(map(str, nums)))   # "9534330"

# -- Stable sort: Python's sort is always stable
# Equal elements maintain original relative order

# -- Sort with multiple criteria (tuples)
students = [("Alice", 90, 20), ("Bob", 85, 22), ("Carol", 90, 19)]
students.sort(key=lambda s: (-s[1], s[2]))  # high grade first, then younger
\end{lstlisting}

%%%==========================================================
\chapter{Searching Algorithms}
%%%==========================================================

\section{Binary Search}

\begin{lstlisting}
# -- Standard Binary Search -- O(log n)
def binary_search(arr, target):
    lo, hi = 0, len(arr) - 1
    while lo <= hi:
        mid = (lo + hi) // 2
        if   arr[mid] == target: return mid
        elif arr[mid] <  target: lo = mid + 1
        else:                    hi = mid - 1
    return -1

# -- First occurrence (leftmost)
def lower_bound(arr, target):
    lo, hi = 0, len(arr)
    while lo < hi:
        mid = (lo + hi) // 2
        if arr[mid] < target: lo = mid + 1
        else:                  hi = mid
    return lo    # == bisect.bisect_left(arr, target)

# -- First position > target (rightmost + 1)
def upper_bound(arr, target):
    lo, hi = 0, len(arr)
    while lo < hi:
        mid = (lo + hi) // 2
        if arr[mid] <= target: lo = mid + 1
        else:                   hi = mid
    return lo    # == bisect.bisect_right(arr, target)

# -- Binary search on answer ("minimise the maximum" / feasibility)
def binary_search_answer(lo, hi, feasible_fn):
    while lo < hi:
        mid = (lo + hi) // 2
        if feasible_fn(mid): hi = mid
        else:                lo = mid + 1
    return lo

# Example: Minimum days to make m bouquets
def min_days_bouquets(bloomDay, m, k):
    if m * k > len(bloomDay): return -1
    def feasible(day):
        bouquets = consec = 0
        for d in bloomDay:
            if d <= day: consec += 1
            else:         consec = 0
            if consec == k: bouquets += 1; consec = 0
        return bouquets >= m
    return binary_search_answer(1, max(bloomDay), feasible)

# -- Find Peak Element -- O(log n)
def find_peak(nums):
    lo, hi = 0, len(nums) - 1
    while lo < hi:
        mid = (lo + hi) // 2
        if nums[mid] < nums[mid+1]: lo = mid + 1
        else:                        hi = mid
    return lo

# -- Search in Rotated Sorted Array -- O(log n)
def search_rotated(nums, target):
    lo, hi = 0, len(nums) - 1
    while lo <= hi:
        mid = (lo + hi) // 2
        if nums[mid] == target: return mid
        if nums[lo] <= nums[mid]:    # left half sorted
            if nums[lo] <= target < nums[mid]: hi = mid - 1
            else:                               lo = mid + 1
        else:                        # right half sorted
            if nums[mid] < target <= nums[hi]: lo = mid + 1
            else:                               hi = mid - 1
    return -1

# -- Find Minimum in Rotated Sorted Array -- O(log n)
def find_min_rotated(nums):
    lo, hi = 0, len(nums) - 1
    while lo < hi:
        mid = (lo + hi) // 2
        if nums[mid] > nums[hi]: lo = mid + 1
        else:                     hi = mid
    return nums[lo]

# -- Binary search on 2D matrix (sorted rows and cols) -- O(m+n)
def search_matrix(matrix, target):
    m, n = len(matrix), len(matrix[0])
    r, c = 0, n - 1           # start from top-right corner
    while r < m and c >= 0:
        if   matrix[r][c] == target: return True
        elif matrix[r][c] <  target: r += 1
        else:                         c -= 1
    return False

# -- Median of Two Sorted Arrays -- O(log(min(m,n)))
def find_median_sorted(nums1, nums2):
    if len(nums1) > len(nums2): nums1, nums2 = nums2, nums1
    m, n   = len(nums1), len(nums2)
    lo, hi = 0, m
    while lo <= hi:
        i  = (lo + hi) // 2
        j  = (m + n + 1) // 2 - i
        max_l1 = float('-inf') if i == 0 else nums1[i-1]
        min_r1 = float('inf')  if i == m else nums1[i]
        max_l2 = float('-inf') if j == 0 else nums2[j-1]
        min_r2 = float('inf')  if j == n else nums2[j]
        if max_l1 <= min_r2 and max_l2 <= min_r1:
            if (m + n) % 2 == 1: return float(max(max_l1, max_l2))
            return (max(max_l1, max_l2) + min(min_r1, min_r2)) / 2.0
        elif max_l1 > min_r2: hi = i - 1
        else:                  lo = i + 1
\end{lstlisting}

\section{Ternary Search}

\begin{lstlisting}
# -- Ternary Search (find maximum of unimodal function) -- O(log n)
def ternary_search(lo, hi, f, eps=1e-9):
    while hi - lo > eps:
        m1 = lo + (hi - lo) / 3
        m2 = hi - (hi - lo) / 3
        if f(m1) < f(m2): lo = m1
        else:              hi = m2
    return (lo + hi) / 2

# Integer version
def ternary_search_int(lo, hi, arr):
    while hi - lo > 2:
        m1 = lo + (hi - lo) // 3
        m2 = hi - (hi - lo) // 3
        if arr[m1] < arr[m2]: lo = m1
        else:                   hi = m2
    return max(range(lo, hi+1), key=lambda i: arr[i])
\end{lstlisting}

\begin{compbox}[Sorting Complexity Summary]
\begin{tabular}{@{}lllll@{}}
\toprule
\textbf{Algorithm} & \textbf{Best} & \textbf{Average} & \textbf{Worst} & \textbf{Stable} \\
\midrule
Bubble Sort    & $O(n)$       & $O(n^2)$      & $O(n^2)$      & Yes \\
Selection Sort & $O(n^2)$     & $O(n^2)$      & $O(n^2)$      & No  \\
Insertion Sort & $O(n)$       & $O(n^2)$      & $O(n^2)$      & Yes \\
Merge Sort     & $O(n\log n)$ & $O(n\log n)$  & $O(n\log n)$  & Yes \\
Quick Sort     & $O(n\log n)$ & $O(n\log n)$  & $O(n^2)$      & No  \\
Heap Sort      & $O(n\log n)$ & $O(n\log n)$  & $O(n\log n)$  & No  \\
Counting Sort  & $O(n+k)$     & $O(n+k)$      & $O(n+k)$      & Yes \\
Radix Sort     & $O(dn)$      & $O(dn)$       & $O(dn)$       & Yes \\
Timsort (Python) & $O(n)$     & $O(n\log n)$  & $O(n\log n)$  & Yes \\
\bottomrule
\end{tabular}
\end{compbox}


%%%%%%%%%%%%%%%%%%%%%%%%%%%%%%%%%%%%%%%%%%%%%%%%%%%%%%%%%%%%%%%
\part{Algorithm Design Paradigms}
%%%%%%%%%%%%%%%%%%%%%%%%%%%%%%%%%%%%%%%%%%%%%%%%%%%%%%%%%%%%%%%

%%%==========================================================
\chapter{Dynamic Programming}
%%%==========================================================

\begin{notebox}[DP in Python vs C++]
Python's \texttt{@cache} / \texttt{@lru\_cache} makes top-down memoisation trivially easy. For bottom-up, use lists. Always call \texttt{sys.setrecursionlimit} for deep recursion.
\end{notebox}

\section{1-D DP}

\begin{lstlisting}
import sys
from functools import cache
sys.setrecursionlimit(300000)

# -- Fibonacci -- O(n) time, O(n) -> O(1) space
@cache
def fib(n):
    if n <= 1: return n
    return fib(n-1) + fib(n-2)

# Bottom-up O(1) space
def fib_iterative(n):
    a, b = 0, 1
    for _ in range(n): a, b = b, a + b
    return a

# -- Climbing Stairs -- O(n)
def climb_stairs(n):
    if n <= 2: return n
    a, b = 1, 2
    for _ in range(3, n+1): a, b = b, a + b
    return b

# -- House Robber -- O(n) time, O(1) space
def rob(nums):
    prev2 = prev1 = 0
    for x in nums:
        prev2, prev1 = prev1, max(prev1, prev2 + x)
    return prev1

# House Robber II (circular array)
def rob2(nums):
    def rob_linear(arr):
        p2 = p1 = 0
        for x in arr: p2, p1 = p1, max(p1, p2 + x)
        return p1
    return max(nums[0], rob_linear(nums[1:]), rob_linear(nums[:-1]))

# -- Coin Change (minimum coins) -- O(n * amount)
def coin_change(coins, amount):
    dp = [float('inf')] * (amount + 1)
    dp[0] = 0
    for c in coins:
        for a in range(c, amount + 1):
            dp[a] = min(dp[a], dp[a - c] + 1)
    return dp[amount] if dp[amount] != float('inf') else -1

# -- Coin Change II (number of ways) -- O(n * amount)
def change(amount, coins):
    dp = [0] * (amount + 1)
    dp[0] = 1
    for c in coins:                  # outer loop on coins -> unbounded knapsack
        for a in range(c, amount+1):
            dp[a] += dp[a - c]
    return dp[amount]

# -- Word Break -- O(n^2 * m) where m = avg word length
def word_break(s, word_dict):
    word_set = set(word_dict)
    n  = len(s)
    dp = [False] * (n + 1)
    dp[0] = True
    for i in range(1, n + 1):
        for j in range(i):
            if dp[j] and s[j:i] in word_set:
                dp[i] = True; break
    return dp[n]

# -- Decode Ways -- O(n)
def num_decodings(s):
    n  = len(s)
    dp = [0] * (n + 1)
    dp[0] = 1
    dp[1] = 0 if s[0] == '0' else 1
    for i in range(2, n + 1):
        one = int(s[i-1:i])
        two = int(s[i-2:i])
        if 1 <= one <= 9:   dp[i] += dp[i-1]
        if 10 <= two <= 26: dp[i] += dp[i-2]
    return dp[n]

# -- Jump Game II (minimum jumps) -- O(n)
def jump(nums):
    jumps = cur_end = cur_far = 0
    for i in range(len(nums) - 1):
        cur_far = max(cur_far, i + nums[i])
        if i == cur_end:
            jumps += 1; cur_end = cur_far
    return jumps
\end{lstlisting}

\section{2-D DP}

\begin{lstlisting}
# -- 0/1 Knapsack -- O(n * W)
def knapsack_01(weights, values, W):
    n  = len(weights)
    dp = [[0]*(W+1) for _ in range(n+1)]
    for i in range(1, n+1):
        for w in range(W+1):
            dp[i][w] = dp[i-1][w]
            if weights[i-1] <= w:
                dp[i][w] = max(dp[i][w], dp[i-1][w-weights[i-1]] + values[i-1])
    return dp[n][W]

# Space optimised O(W)
def knapsack_01_opt(weights, values, W):
    dp = [0] * (W + 1)
    for w_i, v_i in zip(weights, values):
        for w in range(W, w_i - 1, -1):    # traverse backward
            dp[w] = max(dp[w], dp[w - w_i] + v_i)
    return dp[W]

# -- Unbounded Knapsack -- O(n * W)
def knapsack_unbounded(weights, values, W):
    dp = [0] * (W + 1)
    for w in range(1, W + 1):
        for w_i, v_i in zip(weights, values):
            if w_i <= w:
                dp[w] = max(dp[w], dp[w - w_i] + v_i)
    return dp[W]

# -- Longest Common Subsequence -- O(m*n)
def lcs(s, t):
    m, n = len(s), len(t)
    dp = [[0]*(n+1) for _ in range(m+1)]
    for i in range(1, m+1):
        for j in range(1, n+1):
            if s[i-1] == t[j-1]: dp[i][j] = dp[i-1][j-1] + 1
            else:                  dp[i][j] = max(dp[i-1][j], dp[i][j-1])
    return dp[m][n]

# LCS reconstruction
def lcs_string(s, t):
    m, n = len(s), len(t)
    dp = [[0]*(n+1) for _ in range(m+1)]
    for i in range(1, m+1):
        for j in range(1, n+1):
            if s[i-1] == t[j-1]: dp[i][j] = dp[i-1][j-1] + 1
            else:                  dp[i][j] = max(dp[i-1][j], dp[i][j-1])
    result = []
    i, j = m, n
    while i > 0 and j > 0:
        if s[i-1] == t[j-1]: result.append(s[i-1]); i -= 1; j -= 1
        elif dp[i-1][j] > dp[i][j-1]: i -= 1
        else: j -= 1
    return "".join(reversed(result))

# -- Longest Increasing Subsequence -- O(n log n) with patience sort
import bisect
def lis(nums):
    tails = []
    for x in nums:
        pos = bisect.bisect_left(tails, x)
        if pos == len(tails): tails.append(x)
        else:                  tails[pos] = x
    return len(tails)

# LIS O(n^2) with reconstruction
def lis_with_path(nums):
    n = len(nums)
    dp     = [1] * n
    parent = [-1] * n
    for i in range(1, n):
        for j in range(i):
            if nums[j] < nums[i] and dp[j] + 1 > dp[i]:
                dp[i] = dp[j] + 1; parent[i] = j
    end = dp.index(max(dp))
    path = []
    while end != -1: path.append(nums[end]); end = parent[end]
    return max(dp), path[::-1]

# -- Edit Distance (Levenshtein) -- O(m*n)
def edit_distance(s, t):
    m, n = len(s), len(t)
    dp = list(range(n + 1))
    for i in range(1, m+1):
        prev = dp[:]
        dp[0] = i
        for j in range(1, n+1):
            if s[i-1] == t[j-1]: dp[j] = prev[j-1]
            else: dp[j] = 1 + min(prev[j], dp[j-1], prev[j-1])
    return dp[n]

# -- Unique Paths -- O(m*n)
def unique_paths(m, n):
    dp = [1] * n
    for _ in range(1, m):
        for j in range(1, n): dp[j] += dp[j-1]
    return dp[n-1]

# -- Minimum Path Sum -- O(m*n)
def min_path_sum(grid):
    m, n = len(grid), len(grid[0])
    dp = grid[0][:]
    for j in range(1, n): dp[j] += dp[j-1]
    for i in range(1, m):
        dp[0] += grid[i][0]
        for j in range(1, n):
            dp[j] = grid[i][j] + min(dp[j], dp[j-1])
    return dp[n-1]

# -- Regular Expression Matching -- O(m*n)
def is_match(s, p):
    m, n = len(s), len(p)
    dp = [[False]*(n+1) for _ in range(m+1)]
    dp[0][0] = True
    for j in range(1, n+1):
        if p[j-1] == '*': dp[0][j] = dp[0][j-2]
    for i in range(1, m+1):
        for j in range(1, n+1):
            if p[j-1] in (s[i-1], '.'):
                dp[i][j] = dp[i-1][j-1]
            elif p[j-1] == '*':
                dp[i][j] = dp[i][j-2]    # zero occurrences
                if p[j-2] in (s[i-1], '.'):
                    dp[i][j] = dp[i][j] or dp[i-1][j]
    return dp[m][n]
\end{lstlisting}

\section{Interval DP}

\begin{lstlisting}
# -- Matrix Chain Multiplication -- O(n^3)
def matrix_chain(dims):
    n = len(dims) - 1    # number of matrices
    dp = [[0]*n for _ in range(n)]
    for length in range(2, n+1):
        for i in range(n - length + 1):
            j = i + length - 1
            dp[i][j] = float('inf')
            for k in range(i, j):
                cost = dp[i][k] + dp[k+1][j] + dims[i]*dims[k+1]*dims[j+1]
                dp[i][j] = min(dp[i][j], cost)
    return dp[0][n-1]

# -- Burst Balloons -- O(n^3)
def max_coins(nums):
    nums = [1] + nums + [1]
    n = len(nums)
    dp = [[0]*n for _ in range(n)]
    for length in range(2, n):
        for lo in range(n - length):
            hi = lo + length
            for k in range(lo+1, hi):
                dp[lo][hi] = max(dp[lo][hi],
                    dp[lo][k] + nums[lo]*nums[k]*nums[hi] + dp[k][hi])
    return dp[0][n-1]

# -- Minimum Cost to Cut a Stick -- O(n^3)
def min_cost_cut(n, cuts):
    cuts = [0] + sorted(cuts) + [n]
    m = len(cuts)
    dp = [[0]*m for _ in range(m)]
    for length in range(2, m):
        for lo in range(m - length):
            hi = lo + length
            dp[lo][hi] = min(cuts[hi]-cuts[lo] + dp[lo][k] + dp[k][hi]
                             for k in range(lo+1, hi))
    return dp[0][m-1]

# -- Palindrome Partitioning II (min cuts) -- O(n^2)
def min_cut(s):
    n = len(s)
    is_pal = [[False]*n for _ in range(n)]
    for i in range(n-1, -1, -1):
        for j in range(i, n):
            if s[i] == s[j] and (j-i <= 2 or is_pal[i+1][j-1]):
                is_pal[i][j] = True
    dp = list(range(-1, n))    # dp[i] = min cuts for s[0..i]
    for i in range(n):
        for j in range(i+1):
            if is_pal[j][i]:
                dp[i+1] = min(dp[i+1], dp[j] + 1)
    return dp[n]
\end{lstlisting}

\section{DP on Trees}

\begin{lstlisting}
# -- Maximum Independent Set on Tree -- O(n)
def tree_max_independent_set(adj, n):
    dp = [[0, 0] for _ in range(n)]    # [excl_root, incl_root]

    def dfs(u, parent):
        dp[u][1] = 1    # include u
        for v in adj[u]:
            if v == parent: continue
            dfs(v, u)
            dp[u][0] += max(dp[v][0], dp[v][1])
            dp[u][1] += dp[v][0]

    dfs(0, -1)
    return max(dp[0][0], dp[0][1])

# -- Diameter of Tree (longest path) -- O(n)
def tree_diameter(adj, n):
    dist = [0] * n
    def dfs(u, p):
        max1 = max2 = 0
        for v in adj[u]:
            if v == p: continue
            dfs(v, u)
            d = dist[v] + 1
            if d > max1:    max2, max1 = max1, d
            elif d > max2:  max2 = d
        dist[u] = max1
        return max1 + max2    # path through u
    ans = [0]
    def dfs2(u, p):
        r = dfs(u, p)
        ans[0] = max(ans[0], r)
    dfs2(0, -1)
    return ans[0]
\end{lstlisting}

\section{Digit DP}

\begin{lstlisting}
# -- Count numbers in [0, n] with digit sum <= S -- O(digits * S)
def count_digit_sum(n, S):
    digits = [int(d) for d in str(n)]
    from functools import cache

    @cache
    def dp(pos, current_sum, tight, started):
        if current_sum > S: return 0
        if pos == len(digits):
            return 1 if started else 0
        limit = digits[pos] if tight else 9
        result = 0
        for d in range(0, limit + 1):
            result += dp(pos+1,
                         current_sum + (d if started or d > 0 else 0),
                         tight and d == limit,
                         started or d > 0)
        return result

    return dp(0, 0, True, False)
\end{lstlisting}

%%%==========================================================
\chapter{Greedy Algorithms}
%%%==========================================================

\begin{lstlisting}
# -- Activity Selection (max non-overlapping intervals) -- O(n log n)
def activity_selection(intervals):
    intervals.sort(key=lambda x: x[1])    # sort by end time
    count = 0; last_end = float('-inf')
    for start, end in intervals:
        if start >= last_end:
            count += 1; last_end = end
    return count

# -- Minimum Number of Platforms (trains) -- O(n log n)
def min_platforms(arrivals, departures):
    arrivals.sort(); departures.sort()
    platforms = max_platforms = 0
    i = j = 0
    while i < len(arrivals):
        if arrivals[i] <= departures[j]:
            platforms += 1; i += 1
            max_platforms = max(max_platforms, platforms)
        else:
            platforms -= 1; j += 1
    return max_platforms

# -- Gas Station (circular tour) -- O(n)
def can_complete_circuit(gas, cost):
    total = tank = start = 0
    for i in range(len(gas)):
        diff = gas[i] - cost[i]
        total += diff; tank += diff
        if tank < 0: start = i + 1; tank = 0
    return start if total >= 0 else -1

# -- Jump Game (can reach end?) -- O(n)
def can_jump(nums):
    max_reach = 0
    for i, x in enumerate(nums):
        if i > max_reach: return False
        max_reach = max(max_reach, i + x)
    return True

# -- Assign Cookies -- O(n log n)
def find_content_children(greed, size):
    greed.sort(); size.sort()
    child = cookie = 0
    while child < len(greed) and cookie < len(size):
        if size[cookie] >= greed[child]: child += 1
        cookie += 1
    return child

# -- Fractional Knapsack -- O(n log n)
def fractional_knapsack(weights, values, W):
    items = sorted(zip(values, weights), key=lambda x: x[0]/x[1], reverse=True)
    total = 0.0
    for v, w in items:
        if W <= 0: break
        take   = min(w, W)
        total += take * (v / w)
        W     -= take
    return total

# -- Huffman Encoding -- O(n log n)
import heapq
def huffman(freq):
    heap = [[f, [ch, ""]] for ch, f in freq.items()]
    heapq.heapify(heap)
    while len(heap) > 1:
        lo = heapq.heappop(heap)
        hi = heapq.heappop(heap)
        for pair in lo[1:]: pair[1] = '0' + pair[1]
        for pair in hi[1:]: pair[1] = '1' + pair[1]
        heapq.heappush(heap, [lo[0]+hi[0]] + lo[1:] + hi[1:])
    return {pair[0]: pair[1] for pair in heap[0][1:]}

# -- Minimum Spanning Tree cost (Kruskal/Prim) -- see Graph chapter

# -- Largest Number -- O(n log n)
from functools import cmp_to_key
def largest_number(nums):
    nums = list(map(str, nums))
    nums.sort(key=cmp_to_key(lambda a, b: 1 if a+b > b+a else -1), reverse=True)
    return "".join(nums).lstrip('0') or "0"

# -- Candy (minimum candies) -- O(n)
def candy(ratings):
    n = len(ratings)
    candies = [1] * n
    for i in range(1, n):         # left pass
        if ratings[i] > ratings[i-1]: candies[i] = candies[i-1] + 1
    for i in range(n-2, -1, -1):  # right pass
        if ratings[i] > ratings[i+1]: candies[i] = max(candies[i], candies[i+1]+1)
    return sum(candies)

# -- Meeting Rooms II (minimum conference rooms) -- O(n log n)
def min_meeting_rooms(intervals):
    intervals.sort(key=lambda x: x[0])
    heap = []    # min-heap of end times
    for start, end in intervals:
        if heap and heap[0] <= start: heapq.heappop(heap)
        heapq.heappush(heap, end)
    return len(heap)

# -- Minimum Cost to Hire K Workers -- O(n log n)
def mincost_to_hire_workers(quality, wage, k):
    workers = sorted(zip(wage, quality), key=lambda x: x[0]/x[1])
    max_heap = []    # max-heap of quality (negated)
    q_sum = 0; ans = float('inf')
    for w, q in workers:
        heapq.heappush(max_heap, -q); q_sum += q
        if len(max_heap) > k: q_sum += heapq.heappop(max_heap)
        if len(max_heap) == k: ans = min(ans, w/q * q_sum)
    return ans
\end{lstlisting}

%%%==========================================================
\chapter{Divide and Conquer}
%%%==========================================================

\begin{lstlisting}
# -- Merge Sort (canonical D&C) -- O(n log n)
# See Sorting chapter.

# -- Quick Select (kth smallest) -- O(n) avg
def quick_select(nums, k):
    # Returns kth smallest (1-indexed)
    def select(lo, hi, k):
        if lo == hi: return nums[lo]
        pivot_idx = _partition(nums, lo, hi)
        rank = pivot_idx - lo + 1
        if rank == k:   return nums[pivot_idx]
        elif rank > k:  return select(lo, pivot_idx-1, k)
        else:           return select(pivot_idx+1, hi, k-rank)
    return select(0, len(nums)-1, k)

# -- Maximum Subarray (D&C) -- O(n log n)
def max_subarray_dc(arr, lo=None, hi=None):
    if lo is None: lo, hi = 0, len(arr)-1
    if lo == hi: return arr[lo]
    mid = (lo + hi) // 2
    left  = max_subarray_dc(arr, lo, mid)
    right = max_subarray_dc(arr, mid+1, hi)
    # Max crossing subarray
    left_max = curr = 0
    for i in range(mid, lo-1, -1): curr += arr[i]; left_max = max(left_max, curr)
    right_max = curr = 0
    for i in range(mid+1, hi+1):   curr += arr[i]; right_max = max(right_max, curr)
    return max(left, right, left_max + right_max)

# -- Closest Pair of Points -- O(n log n)
def closest_pair(points):
    import math
    def dist(p, q): return math.hypot(p[0]-q[0], p[1]-q[1])

    def closest_strip(strip, d):
        strip.sort(key=lambda p: p[1])
        min_d = d
        for i in range(len(strip)):
            j = i + 1
            while j < len(strip) and strip[j][1] - strip[i][1] < min_d:
                min_d = min(min_d, dist(strip[i], strip[j])); j += 1
        return min_d

    def rec(pts):
        n = len(pts)
        if n <= 3:
            min_d = float('inf')
            for i in range(n):
                for j in range(i+1, n): min_d = min(min_d, dist(pts[i], pts[j]))
            return min_d
        mid  = n // 2
        mid_x = pts[mid][0]
        d = min(rec(pts[:mid]), rec(pts[mid:]))
        strip = [p for p in pts if abs(p[0]-mid_x) < d]
        return min(d, closest_strip(strip, d))

    points.sort()
    return rec(points)

# -- Fast Power (Binary Exponentiation) -- O(log n)
def fast_pow(base, exp, mod=None):
    result = 1
    base = base % mod if mod else base
    while exp > 0:
        if exp & 1:
            result = (result * base) % mod if mod else result * base
        base = (base * base) % mod if mod else base * base
        exp >>= 1
    return result

# -- Karatsuba Multiplication -- O(n^1.585)
def karatsuba(x, y):
    if x < 10 or y < 10: return x * y
    n = max(len(str(x)), len(str(y)))
    m = n // 2
    high_x, low_x = divmod(x, 10**m)
    high_y, low_y = divmod(y, 10**m)
    z0 = karatsuba(low_x, low_y)
    z2 = karatsuba(high_x, high_y)
    z1 = karatsuba(low_x + high_x, low_y + high_y) - z0 - z2
    return z2 * 10**(2*m) + z1 * 10**m + z0
\end{lstlisting}

%%%==========================================================
\chapter{Backtracking}
%%%==========================================================

\begin{lstlisting}
# -- Permutations -- O(n * n!)
def permutations(nums):
    result = []
    def backtrack(path, used):
        if len(path) == len(nums): result.append(path[:]); return
        for i, x in enumerate(nums):
            if used[i]: continue
            used[i] = True; path.append(x)
            backtrack(path, used)
            path.pop(); used[i] = False
    backtrack([], [False]*len(nums))
    return result

# Permutations with duplicates
def permutations_unique(nums):
    result = []; nums.sort()
    def backtrack(path, used):
        if len(path) == len(nums): result.append(path[:]); return
        for i in range(len(nums)):
            if used[i]: continue
            if i > 0 and nums[i] == nums[i-1] and not used[i-1]: continue
            used[i] = True; path.append(nums[i])
            backtrack(path, used)
            path.pop(); used[i] = False
    backtrack([], [False]*len(nums))
    return result

# -- Subsets -- O(2^n)
def subsets(nums):
    result = []
    def backtrack(start, path):
        result.append(path[:])
        for i in range(start, len(nums)):
            path.append(nums[i])
            backtrack(i+1, path)
            path.pop()
    backtrack(0, [])
    return result

# Subsets with duplicates
def subsets_unique(nums):
    result = []; nums.sort()
    def backtrack(start, path):
        result.append(path[:])
        for i in range(start, len(nums)):
            if i > start and nums[i] == nums[i-1]: continue
            path.append(nums[i]); backtrack(i+1, path); path.pop()
    backtrack(0, [])
    return result

# -- Combinations -- O(C(n,k))
def combinations(n, k):
    result = []
    def backtrack(start, path):
        if len(path) == k: result.append(path[:]); return
        for i in range(start, n+1):
            path.append(i); backtrack(i+1, path); path.pop()
    backtrack(1, [])
    return result

# Combination Sum (unlimited use, target sum)
def combination_sum(candidates, target):
    result = []; candidates.sort()
    def backtrack(start, path, remaining):
        if remaining == 0: result.append(path[:]); return
        for i in range(start, len(candidates)):
            if candidates[i] > remaining: break
            path.append(candidates[i])
            backtrack(i, path, remaining - candidates[i])    # i not i+1 (reuse)
            path.pop()
    backtrack(0, [], target)
    return result

# Combination Sum II (each number used once)
def combination_sum2(candidates, target):
    result = []; candidates.sort()
    def backtrack(start, path, remaining):
        if remaining == 0: result.append(path[:]); return
        for i in range(start, len(candidates)):
            if candidates[i] > remaining: break
            if i > start and candidates[i] == candidates[i-1]: continue
            path.append(candidates[i])
            backtrack(i+1, path, remaining - candidates[i])
            path.pop()
    backtrack(0, [], target)
    return result

# -- N-Queens -- O(n!)
def solve_n_queens(n):
    result = []
    cols = set(); diag1 = set(); diag2 = set()   # diagonals

    def backtrack(row, board):
        if row == n: result.append(["".join(r) for r in board]); return
        for col in range(n):
            if col in cols or (row-col) in diag1 or (row+col) in diag2: continue
            cols.add(col); diag1.add(row-col); diag2.add(row+col)
            board[row][col] = 'Q'
            backtrack(row+1, board)
            board[row][col] = '.'; cols.discard(col)
            diag1.discard(row-col); diag2.discard(row+col)

    board = [['.']*n for _ in range(n)]
    backtrack(0, board)
    return result

# -- Sudoku Solver -- O(9^(n*n)) but in practice much faster
def solve_sudoku(board):
    empty = [(r, c) for r in range(9) for c in range(9) if board[r][c] == '.']

    def is_valid(r, c, ch):
        if ch in board[r]: return False
        if ch in [board[i][c] for i in range(9)]: return False
        br, bc = 3*(r//3), 3*(c//3)
        for i in range(br, br+3):
            for j in range(bc, bc+3):
                if board[i][j] == ch: return False
        return True

    def backtrack(idx):
        if idx == len(empty): return True
        r, c = empty[idx]
        for ch in '123456789':
            if is_valid(r, c, ch):
                board[r][c] = ch
                if backtrack(idx+1): return True
                board[r][c] = '.'
        return False

    backtrack(0)

# -- Word Search in Grid -- O(m*n*4^L)
def word_search(board, word):
    m, n = len(board), len(board[0])
    def dfs(r, c, idx):
        if idx == len(word): return True
        if r < 0 or r >= m or c < 0 or c >= n: return False
        if board[r][c] != word[idx]: return False
        tmp = board[r][c]; board[r][c] = '#'
        found = any(dfs(r+dr, c+dc, idx+1) for dr, dc in [(-1,0),(1,0),(0,-1),(0,1)])
        board[r][c] = tmp
        return found
    return any(dfs(r, c, 0) for r in range(m) for c in range(n))

# -- Letter Combinations of Phone Number -- O(4^n)
def letter_combinations(digits):
    if not digits: return []
    mapping = {'2':'abc','3':'def','4':'ghi','5':'jkl',
               '6':'mno','7':'pqrs','8':'tuv','9':'wxyz'}
    result = []
    def backtrack(idx, path):
        if idx == len(digits): result.append("".join(path)); return
        for ch in mapping[digits[idx]]:
            path.append(ch); backtrack(idx+1, path); path.pop()
    backtrack(0, [])
    return result
\end{lstlisting}


%%%%%%%%%%%%%%%%%%%%%%%%%%%%%%%%%%%%%%%%%%%%%%%%%%%%%%%%%%%%%%%
\part{String and Mathematical Algorithms}
%%%%%%%%%%%%%%%%%%%%%%%%%%%%%%%%%%%%%%%%%%%%%%%%%%%%%%%%%%%%%%%

%%%==========================================================
\chapter{String Algorithms}
%%%==========================================================

\section{KMP Pattern Matching}

\begin{lstlisting}
# -- KMP (Knuth-Morris-Pratt) -- O(n + m)
def kmp_search(text, pattern):
    def build_lps(pattern):
        m   = len(pattern)
        lps = [0] * m       # longest proper prefix which is also suffix
        length = 0; i = 1
        while i < m:
            if pattern[i] == pattern[length]:
                length += 1; lps[i] = length; i += 1
            else:
                if length != 0: length = lps[length-1]
                else:           lps[i] = 0; i += 1
        return lps

    n, m = len(text), len(pattern)
    if m == 0: return []
    lps  = build_lps(pattern)
    result = []; i = j = 0
    while i < n:
        if text[i] == pattern[j]:
            i += 1; j += 1
        if j == m:
            result.append(i - j)
            j = lps[j-1]
        elif i < n and text[i] != pattern[j]:
            if j != 0: j = lps[j-1]
            else:      i += 1
    return result    # list of start indices

# -- Z-Function -- O(n)
def z_function(s):
    n = len(s)
    z = [0] * n
    z[0] = n
    l = r = 0
    for i in range(1, n):
        if i < r: z[i] = min(r-i, z[i-l])
        while i + z[i] < n and s[z[i]] == s[i + z[i]]:
            z[i] += 1
        if i + z[i] > r: l, r = i, i + z[i]
    return z

def z_search(text, pattern):
    s = pattern + '#' + text
    z = z_function(s)
    m = len(pattern)
    return [i - m - 1 for i in range(m+1, len(s)) if z[i] == m]
\end{lstlisting}

\section{Rabin-Karp and Rolling Hash}

\begin{lstlisting}
# -- Rabin-Karp -- O(n + m) avg, O(nm) worst
def rabin_karp(text, pattern):
    n, m = len(text), len(pattern)
    if m > n: return []
    MOD  = (1 << 61) - 1    # Mersenne prime
    BASE = 131
    pw   = pow(BASE, m, MOD)

    def h(s): return sum(ord(c) * pow(BASE, i, MOD) for i, c in enumerate(s)) % MOD

    p_hash = h(pattern)
    t_hash = h(text[:m])
    result = []
    if t_hash == p_hash and text[:m] == pattern: result.append(0)

    for i in range(1, n - m + 1):
        t_hash = (t_hash - ord(text[i-1]) + MOD) % MOD
        t_hash = t_hash * BASE % MOD
        t_hash = (t_hash + ord(text[i+m-1]) * pow(BASE, m-1, MOD)) % MOD
        if t_hash == p_hash and text[i:i+m] == pattern:
            result.append(i)
    return result

# -- Polynomial Rolling Hash for substring comparison -- O(n) build, O(1) query
class RollingHash:
    def __init__(self, s, base=131, mod=(1<<61)-1):
        n = len(s)
        self.mod  = mod
        self.pw   = [1] * (n+1)
        self.h    = [0] * (n+1)
        for i in range(n):
            self.h[i+1]  = (self.h[i] * base + ord(s[i])) % mod
            self.pw[i+1] = self.pw[i] * base % mod

    def query(self, l, r):    # hash of s[l..r] inclusive, O(1)
        return (self.h[r+1] - self.h[l] * self.pw[r-l+1]) % self.mod
\end{lstlisting}

\section{Suffix Array and LCP Array}

\begin{lstlisting}
# -- Suffix Array (O(n log^2 n) simple version)
def build_suffix_array(s):
    s  += '$'    # sentinel (lexicographically smallest)
    n  = len(s)
    sa = sorted(range(n), key=lambda i: s[i:])
    return sa[1:]    # exclude sentinel

# -- LCP Array (Kasai's algorithm) -- O(n)
def build_lcp(s, sa):
    n    = len(s)
    rank = [0] * n
    for i, suf in enumerate(sa): rank[suf] = i
    lcp  = [0] * n
    h    = 0
    for i in range(n):
        if rank[i] > 0:
            j = sa[rank[i]-1]
            while i+h < n and j+h < n and s[i+h] == s[j+h]: h += 1
            lcp[rank[i]] = h
            if h > 0: h -= 1
    return lcp

# -- Longest Repeated Substring using SA + LCP -- O(n)
def longest_repeated_substring(s):
    sa  = build_suffix_array(s)
    lcp = build_lcp(s, sa)
    max_lcp = max(lcp)
    idx = lcp.index(max_lcp)
    return s[sa[idx]: sa[idx] + max_lcp]
\end{lstlisting}

\section{Manacher's Algorithm}

\begin{lstlisting}
# -- Manacher -- longest palindromic substring -- O(n)
def manacher(s):
    T   = '#' + '#'.join(s) + '#'
    n   = len(T)
    P   = [0] * n
    C = R = 0
    for i in range(n):
        mirror = 2*C - i
        if i < R: P[i] = min(R-i, P[mirror])
        while i+P[i]+1 < n and i-P[i]-1 >= 0 and T[i+P[i]+1] == T[i-P[i]-1]:
            P[i] += 1
        if i+P[i] > R: C, R = i, i+P[i]
    max_len = max(P)
    center  = P.index(max_len)
    start   = (center - max_len) // 2
    return s[start: start + max_len]

# Count all palindromic substrings using Manacher
def count_palindromes(s):
    T = '#' + '#'.join(s) + '#'
    n = len(T); P = [0]*n; C = R = 0
    for i in range(n):
        mirror = 2*C - i
        if i < R: P[i] = min(R-i, P[mirror])
        while i+P[i]+1 < n and i-P[i]-1 >= 0 and T[i+P[i]+1] == T[i-P[i]-1]:
            P[i] += 1
        if i+P[i] > R: C, R = i, i+P[i]
    return sum((p+1)//2 for p in P)
\end{lstlisting}

\section{Aho-Corasick Multi-Pattern Search}

\begin{lstlisting}
# -- Aho-Corasick -- O(n + m + z), z = total matches
from collections import deque, defaultdict

class AhoCorasick:
    def __init__(self, words):
        self.goto   = [{}]
        self.fail   = [0]
        self.output = [[]]
        self._build(words)

    def _build(self, words):
        for word in words:
            cur = 0
            for ch in word:
                if ch not in self.goto[cur]:
                    self.goto.append({}); self.fail.append(0); self.output.append([])
                    self.goto[cur][ch] = len(self.goto) - 1
                cur = self.goto[cur][ch]
            self.output[cur].append(word)
        # BFS to set failure links
        queue = deque()
        for ch, s in self.goto[0].items():
            queue.append(s)
        while queue:
            r = queue.popleft()
            for ch, s in self.goto[r].items():
                queue.append(s)
                state = self.fail[r]
                while state and ch not in self.goto[state]: state = self.fail[state]
                self.fail[s] = self.goto[state].get(ch, 0)
                if self.fail[s] == s: self.fail[s] = 0
                self.output[s] += self.output[self.fail[s]]

    def search(self, text):
        cur = 0; results = []
        for i, ch in enumerate(text):
            while cur and ch not in self.goto[cur]: cur = self.fail[cur]
            cur = self.goto[cur].get(ch, 0)
            for word in self.output[cur]:
                results.append((i - len(word) + 1, word))
        return results
\end{lstlisting}

%%%==========================================================
\chapter{Mathematical Algorithms}
%%%==========================================================

\section{Number Theory}

\begin{lstlisting}
import math

# -- GCD / LCM -- O(log(min(a,b)))
def gcd(a, b): return math.gcd(a, b)           # or: while b: a,b=b,a%b
def lcm(a, b): return a * b // math.gcd(a, b)

# -- Extended Euclidean Algorithm -- O(log(min(a,b)))
def extended_gcd(a, b):
    # Returns (gcd, x, y) such that a*x + b*y = gcd
    if a == 0: return b, 0, 1
    g, x1, y1 = extended_gcd(b % a, a)
    return g, y1 - (b//a)*x1, x1

# -- Modular Inverse -- O(log m)
def mod_inverse(a, m):
    g, x, _ = extended_gcd(a % m, m)
    if g != 1: raise ValueError("Inverse doesn't exist")
    return (x % m + m) % m

# Python 3.8+ shorthand:
# pow(a, -1, m)

# -- Sieve of Eratosthenes -- O(n log log n)
def sieve(n):
    is_prime = [True] * (n+1)
    is_prime[0] = is_prime[1] = False
    for i in range(2, int(n**0.5)+1):
        if is_prime[i]:
            for j in range(i*i, n+1, i): is_prime[j] = False
    return [i for i in range(2, n+1) if is_prime[i]]

# -- Segmented Sieve -- O(n) for primes up to n
def segmented_sieve(lo, hi):
    small_primes = sieve(int(hi**0.5)+1)
    is_prime = [True] * (hi - lo + 1)
    if lo == 1: is_prime[0] = False
    for p in small_primes:
        start = max(p*p, ((lo + p - 1) // p) * p)
        for j in range(start, hi+1, p): is_prime[j-lo] = False
    return [i+lo for i, v in enumerate(is_prime) if v]

# -- Linear Sieve (Euler's sieve) -- O(n) strict
def linear_sieve(n):
    primes  = []
    is_comp = [False] * (n+1)
    for i in range(2, n+1):
        if not is_comp[i]: primes.append(i)
        for p in primes:
            if i * p > n: break
            is_comp[i*p] = True
            if i % p == 0: break
    return primes

# -- Prime Factorisation -- O(sqrt(n))
def prime_factors(n):
    factors = {}
    d = 2
    while d * d <= n:
        while n % d == 0: factors[d] = factors.get(d,0)+1; n //= d
        d += 1
    if n > 1: factors[n] = factors.get(n,0)+1
    return factors

# -- Euler's Totient Function -- O(sqrt(n)) for single n
def totient(n):
    result = n
    p = 2
    while p*p <= n:
        if n % p == 0:
            while n % p == 0: n //= p
            result -= result // p
        p += 1
    if n > 1: result -= result // n
    return result

# Totient sieve -- O(n log log n) for all values up to n
def totient_sieve(n):
    phi = list(range(n+1))
    for i in range(2, n+1):
        if phi[i] == i:   # i is prime
            for j in range(i, n+1, i): phi[j] -= phi[j] // i
    return phi

# -- Fast Power (Binary Exponentiation) -- O(log n)
def fast_pow(b, e, mod=None):
    r = 1
    b = b % mod if mod else b
    while e:
        if e & 1: r = r*b%mod if mod else r*b
        b = b*b%mod if mod else b*b
        e >>= 1
    return r

# -- Matrix Exponentiation -- O(k^3 log n)
def mat_mul(A, B, mod):
    k = len(A)
    C = [[0]*k for _ in range(k)]
    for i in range(k):
        for j in range(k):
            for l in range(k):
                C[i][j] = (C[i][j] + A[i][l]*B[l][j]) % mod
    return C

def mat_pow(M, n, mod):
    k = len(M)
    R = [[int(i==j) for j in range(k)] for i in range(k)]  # identity
    while n:
        if n & 1: R = mat_mul(R, M, mod)
        M = mat_mul(M, M, mod); n >>= 1
    return R

# Fibonacci with matrix exponentiation -- O(log n)
def fib_matrix(n, mod=10**9+7):
    if n <= 1: return n
    M = [[1,1],[1,0]]
    R = mat_pow(M, n-1, mod)
    return R[0][0]

# -- Chinese Remainder Theorem (CRT) -- O(log(min(m1,m2)))
def crt(r1, m1, r2, m2):
    # Find x such that x = r1 (mod m1), x = r2 (mod m2)
    g, p, q = extended_gcd(m1, m2)
    if (r2 - r1) % g != 0: return None, None   # no solution
    lcm_ = m1 * m2 // g
    x = (r1 + m1 * ((r2-r1)//g * p % (m2//g))) % lcm_
    return x, lcm_

# -- Lucas's Theorem -- n choose r mod prime p
def lucas_ncr(n, r, p):
    if r == 0: return 1
    ni, ri = n % p, r % p
    if ri > ni: return 0
    return ncr(ni, ri, p) * lucas_ncr(n//p, r//p, p) % p

# Precompute factorials for nCr mod p
def precompute_factorials(n, mod):
    fact  = [1]*(n+1)
    inv_f = [1]*(n+1)
    for i in range(1, n+1): fact[i] = fact[i-1]*i % mod
    inv_f[n] = pow(fact[n], mod-2, mod)
    for i in range(n-1, -1, -1): inv_f[i] = inv_f[i+1]*(i+1) % mod
    def ncr(n, r):
        if r < 0 or r > n: return 0
        return fact[n] * inv_f[r] % mod * inv_f[n-r] % mod
    return ncr
\end{lstlisting}

\section{Bit Manipulation}

\begin{lstlisting}
# -- Bit operations in Python ──────────────────────────────────
x = 42                          # 0b101010
x & y                           # AND
x | y                           # OR
x ^ y                           # XOR
~x                              # NOT (= -(x+1) in Python!)
x << k                          # left shift by k
x >> k                          # right shift by k (arithmetic)

# Python integers are arbitrary precision -- no overflow!

# -- Common tricks ─────────────────────────────────────────────
x & 1                           # check if odd
x & (x-1)                      # clear lowest set bit; 0 iff power of 2
x & (-x)                       # isolate lowest set bit
x | (x-1)                      # set all bits below lowest set bit
x.bit_length()                  # number of bits (= floor(log2(x)) + 1)
bin(x).count('1')               # popcount (number of 1-bits)
int.bit_count(x)                # Python 3.10+ native popcount

# Check if power of 2
def is_power_of_two(n): return n > 0 and (n & (n-1)) == 0

# Next power of 2
def next_pow2(n):
    p = 1
    while p < n: p <<= 1
    return p

# -- XOR tricks ────────────────────────────────────────────────
# Single number (one element appears odd times)
def single_number(nums):
    result = 0
    for x in nums: result ^= x
    return result

# Single number II (all others appear 3x) -- O(n)
def single_number_3x(nums):
    ones = twos = 0
    for x in nums:
        ones = (ones ^ x) & ~twos
        twos = (twos ^ x) & ~ones
    return ones

# Two single numbers (two elements appear once)
def single_number_two(nums):
    xor = 0
    for x in nums: xor ^= x
    diff_bit = xor & (-xor)    # rightmost differing bit
    a = b = 0
    for x in nums:
        if x & diff_bit: a ^= x
        else:            b ^= x
    return a, b

# -- Bit DP: bitmask DP -- O(2^n * n)
# Travelling Salesman (held-karp)
def tsp(dist, n):
    INF = float('inf')
    dp  = [[INF]*(1<<n) for _ in range(n)]
    dp[0][1] = 0    # start at city 0, visited = {0}
    for mask in range(1, 1<<n):
        for u in range(n):
            if not (mask >> u & 1): continue
            if dp[u][mask] == INF: continue
            for v in range(n):
                if mask >> v & 1: continue
                new_mask = mask | (1<<v)
                dp[v][new_mask] = min(dp[v][new_mask], dp[u][mask] + dist[u][v])
    full = (1<<n) - 1
    return min(dp[u][full] + dist[u][0] for u in range(1, n))

# Count subsets with given XOR
def count_subsets_xor(arr, x):
    from functools import cache
    n = len(arr)
    @cache
    def dp(i, cur_xor):
        if i == n: return 1 if cur_xor == x else 0
        return dp(i+1, cur_xor^arr[i]) + dp(i+1, cur_xor)
    return dp(0, 0)

# -- Subset Enumeration using bitmask -- O(2^n)
def all_subsets_bitmask(arr):
    n = len(arr)
    for mask in range(1 << n):
        subset = [arr[i] for i in range(n) if mask >> i & 1]
        print(subset)

# Enumerate all subsets of a given mask (for SOS DP)
def enumerate_submasks(mask):
    sub = mask
    while sub:
        yield sub
        sub = (sub-1) & mask

# -- Sum over Subsets (SOS) DP -- O(2^n * n)
def sos_dp(arr):
    n = max(arr).bit_length()
    dp = [0] * (1 << n)
    for x in arr: dp[x] += 1
    for i in range(n):
        for mask in range(1 << n):
            if mask >> i & 1:
                dp[mask] += dp[mask ^ (1 << i)]
    return dp    # dp[mask] = sum of arr[x] for all x subset of mask
\end{lstlisting}

\section{Geometry}

\begin{lstlisting}
# -- 2D Point / Vector operations -- O(1)
def cross(o, a, b):
    # Cross product of vectors OA and OB
    # > 0: counterclockwise, < 0: clockwise, = 0: collinear
    return (a[0]-o[0])*(b[1]-o[1]) - (a[1]-o[1])*(b[0]-o[0])

def dot(a, b):  return a[0]*b[0] + a[1]*b[1]
def dist2(a, b): return (a[0]-b[0])**2 + (a[1]-b[1])**2

# -- Convex Hull (Graham Scan) -- O(n log n)
def convex_hull(points):
    points = sorted(set(points))
    if len(points) <= 1: return points
    # Build lower hull
    lower = []
    for p in points:
        while len(lower) >= 2 and cross(lower[-2], lower[-1], p) <= 0:
            lower.pop()
        lower.append(p)
    # Build upper hull
    upper = []
    for p in reversed(points):
        while len(upper) >= 2 and cross(upper[-2], upper[-1], p) <= 0:
            upper.pop()
        upper.append(p)
    return lower[:-1] + upper[:-1]

# -- Point in Polygon (Ray Casting) -- O(n)
def point_in_polygon(point, polygon):
    x, y = point; inside = False
    n = len(polygon); j = n - 1
    for i in range(n):
        xi, yi = polygon[i]; xj, yj = polygon[j]
        if ((yi > y) != (yj > y)) and (x < (xj-xi)*(y-yi)/(yj-yi)+xi):
            inside = not inside
        j = i
    return inside

# -- Line Segment Intersection -- O(1)
def segments_intersect(p1, p2, p3, p4):
    d1 = cross(p3, p4, p1); d2 = cross(p3, p4, p2)
    d3 = cross(p1, p2, p3); d4 = cross(p1, p2, p4)
    if ((d1 > 0 and d2 < 0) or (d1 < 0 and d2 > 0)) and \
       ((d3 > 0 and d4 < 0) or (d3 < 0 and d4 > 0)):
        return True
    # Check collinear cases (omitted for brevity)
    return False
\end{lstlisting}


%%%%%%%%%%%%%%%%%%%%%%%%%%%%%%%%%%%%%%%%%%%%%%%%%%%%%%%%%%%%%%%
\part{Advanced Topics and Competitive Programming}
%%%%%%%%%%%%%%%%%%%%%%%%%%%%%%%%%%%%%%%%%%%%%%%%%%%%%%%%%%%%%%%

%%%==========================================================
\chapter{Competitive Programming Patterns}
%%%==========================================================

\section{Two Heaps Pattern}

\begin{lstlisting}
# -- Find Median from Data Stream -- O(log n) per add, O(1) query
import heapq
class MedianFinder:
    def __init__(self):
        self.lo = []    # max-heap (negated) for lower half
        self.hi = []    # min-heap for upper half

    def add_num(self, num):
        heapq.heappush(self.lo, -num)
        # Ensure lo's max <= hi's min
        if self.hi and -self.lo[0] > self.hi[0]:
            heapq.heappush(self.hi, -heapq.heappop(self.lo))
        # Balance sizes: lo can have at most 1 extra
        if len(self.lo) > len(self.hi) + 1:
            heapq.heappush(self.hi, -heapq.heappop(self.lo))
        elif len(self.hi) > len(self.lo):
            heapq.heappush(self.lo, -heapq.heappop(self.hi))

    def find_median(self):
        if len(self.lo) > len(self.hi): return float(-self.lo[0])
        return (-self.lo[0] + self.hi[0]) / 2.0
\end{lstlisting}

\section{Intervals}

\begin{lstlisting}
# -- Merge Overlapping Intervals -- O(n log n)
def merge_intervals(intervals):
    intervals.sort(); result = [intervals[0]]
    for start, end in intervals[1:]:
        if start <= result[-1][1]: result[-1][1] = max(result[-1][1], end)
        else:                       result.append([start, end])
    return result

# -- Insert Interval -- O(n)
def insert_interval(intervals, new):
    result = []; i = 0; n = len(intervals)
    # Add all intervals ending before new starts
    while i < n and intervals[i][1] < new[0]: result.append(intervals[i]); i += 1
    # Merge overlapping
    while i < n and intervals[i][0] <= new[1]:
        new[0] = min(new[0], intervals[i][0])
        new[1] = max(new[1], intervals[i][1]); i += 1
    result.append(new)
    result.extend(intervals[i:])
    return result

# -- Non-overlapping Intervals (min removals) -- O(n log n)
def erase_overlap_intervals(intervals):
    intervals.sort(key=lambda x: x[1])
    last_end = float('-inf'); removed = 0
    for start, end in intervals:
        if start >= last_end: last_end = end
        else:                  removed += 1
    return removed
\end{lstlisting}

\section{Monotonic Queue / Stack Patterns}

\begin{lstlisting}
# -- Largest Rectangle in Histogram (revisited) -- O(n)
def largest_rectangle_histogram(heights):
    heights = heights + [0]; stack = [-1]; max_area = 0
    for i, h in enumerate(heights):
        while stack[-1] != -1 and heights[stack[-1]] >= h:
            height = heights[stack.pop()]
            width  = i - stack[-1] - 1
            max_area = max(max_area, height * width)
        stack.append(i)
    return max_area

# -- Jump Game VI (max score with window k) -- O(n log k) with heap
def max_result(nums, k):
    from collections import deque
    n  = len(nums)
    dq = deque([0])   # indices, decreasing dp value
    dp = [0] * n; dp[0] = nums[0]
    for i in range(1, n):
        while dq and dq[0] < i - k: dq.popleft()
        dp[i] = dp[dq[0]] + nums[i]
        while dq and dp[dq[-1]] <= dp[i]: dq.pop()
        dq.append(i)
    return dp[n-1]
\end{lstlisting}

\section{Graph Colouring and Matching}

\begin{lstlisting}
# -- Hungarian Algorithm (assignment problem) -- O(n^3)
def hungarian(cost):
    n = len(cost)
    u  = [0]*(n+1); v = [0]*(n+1)
    p  = [0]*(n+1); way = [0]*(n+1)
    for i in range(1, n+1):
        p[0] = i; j0 = 0
        minv = [float('inf')]*(n+1)
        used = [False]*(n+1)
        while True:
            used[j0] = True; i0 = p[j0]; delta = float('inf'); j1 = -1
            for j in range(1, n+1):
                if not used[j]:
                    cur = cost[i0-1][j-1] - u[i0] - v[j]
                    if cur < minv[j]: minv[j] = cur; way[j] = j0
                    if minv[j] < delta: delta = minv[j]; j1 = j
            for j in range(n+1):
                if used[j]: u[p[j]] += delta; v[j] -= delta
                else:        minv[j] -= delta
            j0 = j1
            if p[j0] == 0: break
        while j0:
            p[j0] = p[way[j0]]; j0 = way[j0]
    ans = [-u[0]]   # total cost
    assignment = [0]*n
    for j in range(1, n+1):
        if p[j]: assignment[p[j]-1] = j-1
    return -u[0], assignment
\end{lstlisting}

%%%==========================================================
\chapter{Miscellaneous Algorithms and Tricks}
%%%==========================================================

\section{Reservoir Sampling}

\begin{lstlisting}
# -- Reservoir Sampling: pick k items uniformly at random -- O(n)
import random
def reservoir_sample(stream, k):
    reservoir = []
    for i, item in enumerate(stream):
        if i < k:
            reservoir.append(item)
        else:
            j = random.randint(0, i)
            if j < k: reservoir[j] = item
    return reservoir
\end{lstlisting}

\section{Fisher-Yates Shuffle}

\begin{lstlisting}
# -- Fisher-Yates Shuffle -- O(n), uniform distribution
def shuffle(arr):
    n = len(arr)
    for i in range(n-1, 0, -1):
        j = random.randint(0, i)
        arr[i], arr[j] = arr[j], arr[i]
    return arr
\end{lstlisting}

\section{Meeting in the Middle}

\begin{lstlisting}
# -- 4-sum count (meet-in-the-middle) -- O(n^2)
def four_sum_count(A, B, C, D):
    from collections import Counter
    ab_sums = Counter(a+b for a in A for b in B)
    return sum(ab_sums.get(-(c+d), 0) for c in C for d in D)
\end{lstlisting}

\section{Mo's Algorithm}

\begin{lstlisting}
# -- Mo's Algorithm: offline range queries -- O((n+q)*sqrt(n))
import math
def mo_algorithm(arr, queries):
    n = len(arr); block = max(1, int(math.sqrt(n)))
    # Sort queries: (l, r, idx) by block of l, then r
    queries = sorted(enumerate(queries), key=lambda e: (e[1][0]//block, e[1][1]))

    cur_l = cur_r = 0
    cur_sum = arr[0]
    results = [0] * len(queries)

    def add(pos): nonlocal cur_sum; cur_sum += arr[pos]
    def remove(pos): nonlocal cur_sum; cur_sum -= arr[pos]

    for idx, (l, r) in queries:
        while cur_r < r: cur_r += 1; add(cur_r)
        while cur_l > l: cur_l -= 1; add(cur_l)
        while cur_r > r: remove(cur_r); cur_r -= 1
        while cur_l < l: remove(cur_l); cur_l += 1
        results[idx] = cur_sum

    return results
\end{lstlisting}

\section{Heavy-Light Decomposition (HLD)}

\begin{lstlisting}
# -- Heavy-Light Decomposition -- O(n) build, O(log^2 n) path query
class HLD:
    def __init__(self, adj, values, root=0):
        n = self.n = len(adj)
        self.parent = [-1]*n; self.depth = [0]*n
        self.heavy  = [-1]*n; self.head  = [0]*n
        self.pos    = [0]*n; self.sz     = [1]*n
        self._dfs_sz(root, -1, adj)
        self._dfs_hld(root, root, adj, [0])
        # Build segment tree on HLD positions
        arr = [values[i] for i in range(n)]   # reorder by pos
        self.seg = SegTree(arr)

    def _dfs_sz(self, u, p, adj):
        self.parent[u] = p; heavy_sz = 0
        for v in adj[u]:
            if v == p: continue
            self.depth[v] = self.depth[u]+1
            self._dfs_sz(v, u, adj)
            self.sz[u] += self.sz[v]
            if self.sz[v] > heavy_sz: heavy_sz = self.sz[v]; self.heavy[u] = v

    def _dfs_hld(self, u, h, adj, cur):
        self.head[u] = h; self.pos[u] = cur[0]; cur[0] += 1
        if self.heavy[u] != -1: self._dfs_hld(self.heavy[u], h, adj, cur)
        for v in adj[u]:
            if v != self.parent[u] and v != self.heavy[u]:
                self._dfs_hld(v, v, adj, cur)

    def path_query(self, u, v):
        result = 0
        while self.head[u] != self.head[v]:
            if self.depth[self.head[u]] < self.depth[self.head[v]]: u, v = v, u
            result += self.seg.query(self.pos[self.head[u]], self.pos[u])
            u = self.parent[self.head[u]]
        if self.depth[u] > self.depth[v]: u, v = v, u
        result += self.seg.query(self.pos[u], self.pos[v])
        return result
\end{lstlisting}

\section{Centroid Decomposition}

\begin{lstlisting}
# -- Centroid Decomposition -- O(n log n) build, O(log n) per query
class CentroidDecomp:
    def __init__(self, adj, n):
        self.adj    = adj; self.n = n
        self.sz     = [0]*n; self.removed = [False]*n
        self.parent = [-1]*n

    def _get_sz(self, u, p):
        self.sz[u] = 1
        for v in self.adj[u]:
            if v != p and not self.removed[v]:
                self._get_sz(v, u); self.sz[u] += self.sz[v]

    def _get_centroid(self, u, p, tree_sz):
        for v in self.adj[u]:
            if v != p and not self.removed[v]:
                if self.sz[v] > tree_sz // 2:
                    return self._get_centroid(v, u, tree_sz)
        return u

    def build(self, u=0, p=-1):
        self._get_sz(u, -1)
        c = self._get_centroid(u, -1, self.sz[u])
        self.removed[c] = True
        self.parent[c]  = p
        for v in self.adj[c]:
            if not self.removed[v]: self.build(v, c)
        self.removed[c] = False
        return c
\end{lstlisting}

%%%==========================================================
\chapter{Input / Output Optimisation}
%%%==========================================================

\begin{lstlisting}
import sys
input = sys.stdin.readline    # 3-5x faster than built-in input()

# -- Fast read of integers
def read_int():     return int(sys.stdin.readline())
def read_ints():    return list(map(int, sys.stdin.readline().split()))
def read_matrix(m): return [list(map(int, sys.stdin.readline().split())) for _ in range(m)]

# -- Fast output (buffer before printing)
out = []
out.append(str(answer))
print('\n'.join(out))         # single system call vs n print() calls

# -- Or use sys.stdout.write
sys.stdout.write('\n'.join(map(str, results)) + '\n')

# -- PyPy tips (if available):
# 1. Use sys.stdin.read() and split once upfront
# 2. Avoid excessive function calls
# 3. Use local variables in tight loops (faster than global)

# -- Complete fast I/O template
import sys
from collections import defaultdict, deque
from heapq import *
from bisect import *
from itertools import *
from functools import cache
input = sys.stdin.readline
MOD   = 10**9 + 7
INF   = float('inf')

def solve():
    n = int(input())
    arr = list(map(int, input().split()))
    # ... solution ...

T = int(input())
for _ in range(T):
    solve()
\end{lstlisting}

%%%==========================================================
\chapter{Complete Algorithm Reference Table}
%%%==========================================================

\begin{longtable}{@{}p{4.5cm}p{2.5cm}p{2.5cm}p{3.5cm}@{}}
\toprule
\textbf{Algorithm / DS} & \textbf{Time} & \textbf{Space} & \textbf{Notes} \\
\midrule
\endhead
\multicolumn{4}{@{}l@{}}{\textbf{Linear Structures}} \\
Array access       & $O(1)$        & ---           & Random access \\
Array search       & $O(n)$        & ---           & Linear scan \\
Dynamic array push & $O(1)$ amort  & $O(n)$        & Python list \\
Linked list insert & $O(1)$        & $O(n)$        & With pointer \\
Stack push/pop     & $O(1)$        & $O(n)$        & \\
Queue enqueue/deq  & $O(1)$        & $O(n)$        & deque \\
\midrule
\multicolumn{4}{@{}l@{}}{\textbf{Trees}} \\
BST search/insert  & $O(\log n)$ avg & $O(n)$      & $O(n)$ worst \\
AVL search/insert  & $O(\log n)$   & $O(n)$        & Guaranteed \\
Heap push/pop      & $O(\log n)$   & $O(n)$        & \\
Heap build         & $O(n)$        & $O(n)$        & \\
Trie insert/search & $O(m)$        & $O(nm)$       & $m$ = key length \\
Segment tree query & $O(\log n)$   & $O(4n)$       & \\
BIT update/query   & $O(\log n)$   & $O(n)$        & \\
Sparse table query & $O(1)$        & $O(n\log n)$  & Static only \\
\midrule
\multicolumn{4}{@{}l@{}}{\textbf{Hashing}} \\
Dict get/set/del   & $O(1)$ avg    & $O(n)$        & $O(n)$ worst \\
Set add/in/remove  & $O(1)$ avg    & $O(n)$        & \\
\midrule
\multicolumn{4}{@{}l@{}}{\textbf{Sorting}} \\
Comparison sort    & $O(n\log n)$  & $O(\log n)$   & Lower bound \\
Counting sort      & $O(n+k)$      & $O(k)$        & Integer keys \\
Radix sort         & $O(dn)$       & $O(n+b)$      & $d$ digits, base $b$ \\
\midrule
\multicolumn{4}{@{}l@{}}{\textbf{Graph}} \\
DFS / BFS          & $O(V+E)$      & $O(V)$        & \\
Dijkstra           & $O((V+E)\log V)$ & $O(V)$     & Non-neg weights \\
Bellman-Ford       & $O(VE)$       & $O(V)$        & Neg weights \\
Floyd-Warshall     & $O(V^3)$      & $O(V^2)$      & All-pairs \\
Kruskal MST        & $O(E\log E)$  & $O(V)$        & + DSU \\
Prim MST           & $O((V+E)\log V)$ & $O(V)$     & \\
Topo Sort          & $O(V+E)$      & $O(V)$        & DAG only \\
Tarjan SCC         & $O(V+E)$      & $O(V)$        & \\
Bridges/APs        & $O(V+E)$      & $O(V)$        & \\
\midrule
\multicolumn{4}{@{}l@{}}{\textbf{Strings}} \\
KMP                & $O(n+m)$      & $O(m)$        & \\
Z-function         & $O(n)$        & $O(n)$        & \\
Rabin-Karp         & $O(n+m)$ avg  & $O(1)$        & \\
Manacher           & $O(n)$        & $O(n)$        & Palindromes \\
Suffix Array       & $O(n\log^2 n)$& $O(n)$        & \\
Aho-Corasick       & $O(n+m+z)$    & $O(nm)$       & $z$ = matches \\
\midrule
\multicolumn{4}{@{}l@{}}{\textbf{DP}} \\
1D DP              & $O(n)$--$O(n^2)$ & $O(n)$     & \\
2D DP (LCS, etc.)  & $O(mn)$       & $O(\min(m,n))$ & Optimised \\
LIS                & $O(n\log n)$  & $O(n)$        & Patience sort \\
Knapsack 0/1       & $O(nW)$       & $O(W)$        & \\
Interval DP        & $O(n^3)$      & $O(n^2)$      & \\
Bitmask DP         & $O(2^n \cdot n)$ & $O(2^n)$   & \\
Digit DP           & $O(D \cdot S)$ & $O(D \cdot S)$ & $D$ digits \\
\midrule
\multicolumn{4}{@{}l@{}}{\textbf{Number Theory}} \\
GCD                & $O(\log n)$   & $O(1)$        & \\
Sieve of Eratosthenes & $O(n\log\log n)$ & $O(n)$ & \\
Fast power         & $O(\log n)$   & $O(1)$        & \\
Matrix power       & $O(k^3\log n)$& $O(k^2)$      & $k\times k$ matrix \\
\midrule
\multicolumn{4}{@{}l@{}}{\textbf{Special}} \\
DSU find/union     & $O(\alpha(n))$& $O(n)$        & Nearly $O(1)$ \\
Mo's Algorithm     & $O((n+q)\sqrt{n})$ & $O(n)$  & Offline queries \\
\bottomrule
\end{longtable}

%%%==========================================================
\chapter{Common Pitfalls: C\texttt{++} to Python}
%%%==========================================================

\begin{warnbox}[Top 10 Pitfalls when Moving from C\texttt{++} to Python DSA]
\begin{enumerate}[leftmargin=*]
\item \textbf{Recursion limit:} Python default is 1000. Add \texttt{sys.setrecursionlimit(300000)} or convert to iterative.
\item \textbf{Max-heap:} Python only has min-heap. Negate values: \texttt{heapq.heappush(h, -x)}.
\item \textbf{Integer type:} Python integers never overflow (arbitrary precision). No need for \texttt{long long}.
\item \textbf{Mutable default args:} \texttt{def f(arr=[]):} shares the list across calls. Use \texttt{def f(arr=None): if arr is None: arr = []}
\item \textbf{2D list initialisation:} \texttt{[[0]*n]*m} creates \emph{m} references to the same row. Always use \texttt{[[0]*n for \_ in range(m)]}.
\item \textbf{String concatenation:} \texttt{s += char} in a loop is $O(n^2)$. Use \texttt{"".join(list)}.
\item \textbf{Slow \texttt{list.pop(0)}:} $O(n)$. Use \texttt{collections.deque} for $O(1)$ front pops.
\item \textbf{No built-in balanced BST:} Use \texttt{sortedcontainers.SortedList} (pip install) or simulate with binary search on a list.
\item \textbf{Graph DFS with large input:} Use iterative DFS with explicit stack to avoid recursion limit hits.
\item \textbf{\texttt{is} vs \texttt{==}:} \texttt{is} checks identity, \texttt{==} checks value. For node comparisons use \texttt{is}; for values use \texttt{==}.
\end{enumerate}
\end{warnbox}

\begin{tipbox}[Python Speed-Up Tips for Competitive Programming]
\begin{enumerate}[leftmargin=*]
\item Use \texttt{sys.stdin.readline} instead of \texttt{input()} for fast I/O.
\item Prefer local variables over global in tight loops (Python looks up locals faster).
\item Use \texttt{@cache} or \texttt{@lru\_cache} for memoisation instead of manual dict.
\item \texttt{heapq.heapify} is $O(n)$; don't push one-by-one when building from scratch.
\item \texttt{bisect.bisect\_left/right} for $O(\log n)$ search on sorted lists.
\item Use \texttt{collections.Counter} for frequency maps --- faster than manual dicts.
\item Avoid \texttt{list.sort()} on already-nearly-sorted data with large $k$-inversions.
\item Use \texttt{frozenset} as dict key when you need a set as a hash key.
\item For set-like ordered operations, use \texttt{sortedcontainers.SortedList}.
\item Use \texttt{math.comb(n, r)} for exact nCr (no float rounding).
\end{enumerate}
\end{tipbox}

\end{document}